\documentclass[xcolor=svgnames,9pt,english,aspectratio=169,t]{beamer}
%\documentclass[xcolor=svgnames,9pt,english,handout,aspectratio=169,t]{beamer}
\newcommand{\handout}{0}

% =====================================================================
% POLI 450A (2026) --- shared TEMPLATE layer.
%
% Input from each lecture near the top of its file, BEFORE
% ../course_theme.tex (metropolis must be selected before the course
% palette overrides it):
%
%   \documentclass[xcolor=svgnames,9pt,english,aspectratio=169]{beamer}
%   \newcommand{\handout}{0}
%   % =====================================================================
% POLI 450A (2026) --- shared TEMPLATE layer.
%
% Input from each lecture near the top of its file, BEFORE
% ../course_theme.tex (metropolis must be selected before the course
% palette overrides it):
%
%   \documentclass[xcolor=svgnames,9pt,english,aspectratio=169]{beamer}
%   \newcommand{\handout}{0}
%   % =====================================================================
% POLI 450A (2026) --- shared TEMPLATE layer.
%
% Input from each lecture near the top of its file, BEFORE
% ../course_theme.tex (metropolis must be selected before the course
% palette overrides it):
%
%   \documentclass[xcolor=svgnames,9pt,english,aspectratio=169]{beamer}
%   \newcommand{\handout}{0}
%   \input{../course_preamble.tex}
%   \input{../course_theme.tex}
%   \graphicspath{{./}{./figs/}}
%   \title[Short title]{450A Political Methodology I}
%   \subtitle{Lecture N: ...}
%
% Modeled on the 2026 Northwestern workshop template
% (Teaching/PanelData/_latex/2026/202608_Northwestern/workshop_preamble.tex).
%
% SCOPE: theme, beamer layout, packages, the roadmap hooks, and --- unlike
% the Northwestern decks --- the notation macros. The 450A lectures
% grew from one common ancestor, so their macro blocks agree wherever they
% overlap; merging them is safe. Three points where they had drifted are
% resolved here, and each is flagged inline below.
% =====================================================================

\usetheme{metropolis}
\metroset{sectionpage=none, subsectionpage=none, numbering=none}
\setbeamertemplate{navigation symbols}{}

\setbeamersize{text margin left=5mm, text margin right=5mm}
\setbeamertemplate{itemize/enumerate body begin}{\leftmargini=1.5em\relax}
\setbeamertemplate{itemize/enumerate subbody begin}{\leftmarginii=1.5em\relax}
\setbeamertemplate{itemize/enumerate subsubbody begin}{\leftmarginiii=1.5em\relax}

% ----------------------------------------------------------- top alignment
% Every deck passes the t class option, so beamer applies the frame key
% below to every frame (explicit [plain,c] frames — title page, roadmaps —
% stay centered). Beamer's stock t key puts a rigid .2cm pad above the
% content; frames sized to this template's full-height ceilings then run
% Overfull \vbox by up to that pad. Same key as beamerbaseframe.sty, with
% one change: the pad may shrink to zero when the frame is full.
\makeatletter
\define@key{beamerframe}{t}[true]{% top
  \beamer@frametopskip=.2cm plus .5\paperheight minus .2cm\relax%
  \beamer@framebottomskip=0pt plus 1fill\relax%
  \beamer@frametopskipautobreak=\beamer@frametopskip\relax%
  \beamer@framebottomskipautobreak=\beamer@framebottomskip\relax%
}

% ------------------------------------------------- slide-top check (2026-10-11)
% A slide whose content is too tall for the frame is squeezed: the pad above
% shrinks and its content starts higher than on other slides, so a figure
% "wobbles" when the class steps from one slide to the next (Lecture 4-1,
% slides 1-2, 2026-10-11). The marker below records where each slide's
% content starts (a whatsit: it adds no space and moves nothing).
% `python3 ../check_slides.py lec<N>` then lists every slide that starts
% higher than the deck's normal position. That list must be empty.
% Recorded per slide: page, frame number, content-top y (sp), the room left
% below the title bar, and P for a plain frame (title page, roadmaps).
\providecommand\slidetop[5]{}
\AtBeginDocument{\immediate\write\@auxout{\string\providecommand\string\slidetop[5]{}}}
\def\beamer@markframetop{%
  \pdfsavepos
  \edef\beamer@marktmp{\noexpand\write\noexpand\@auxout{\string\slidetop
    {\noexpand\thepage}{\noexpand\insertframenumber}{\noexpand\the\noexpand\pdflastypos}%
    {\number\beamer@frametextheight}{\ifbeamer@plainframe P\else N\fi}}}%
  \beamer@marktmp}
\RequirePackage{etoolbox}
\patchcmd{\endbeamer@frameslide}
  {\vskip\beamer@frametopskip}
  {\vskip\beamer@frametopskip\beamer@markframetop}
  {}{\PackageWarning{course_preamble}{slide-top marker not installed}}
\makeatother

% ---------------------------------------------------------------- packages
\usepackage{xcolor}
\usepackage{colortbl}
\usepackage{amsmath}
\usepackage{amsfonts}
\usepackage{amssymb}
\usepackage{bm}
\usepackage{graphicx}
\usepackage{multirow}
\usepackage{fancyvrb}
\usepackage{listings}
\usepackage{tikz}
\usetikzlibrary{arrows.meta,positioning,decorations.pathreplacing}

% --------------------------------------------------------- figure fitting
% Default bounding box for every \includegraphics that does not override it.
% With keepaspectratio the image is scaled to fit *inside* width x height, so
% a figure can never run past the side margins or off the bottom of a slide.
% Options given at the call site override these on a per-key basis: e.g.
% [width=.6\linewidth] narrows the box but keeps the height ceiling.
% \figmaxheight is set at \begin{document}: beamer only computes the final
% \textheight once the head/foot templates are known. 0.80 is just under what a
% frame with a title and one line of text can hold (measured: 0.834\textheight);
% a frame that needs more room for text can lower it locally.
\newlength{\figmaxheight}
\AtBeginDocument{\setlength{\figmaxheight}{0.80\textheight}}
\setkeys{Gin}{width=\linewidth,height=\figmaxheight,keepaspectratio}

% -------------------------------------------------------- roadmap frames
% Identical in all lectures. A lecture that does not want the
% subsection roadmap clears it after the \input with \AtBeginSubsection[]{}.
% hideothersubsections is load-bearing, not cosmetic: lecture 4 had 13
% subsections before it was split into 4-1, 4-2 and 4-3 (2026-10-09), and a
% full contents list ran off the top and bottom of the slide. Hiding the other sections' subsections caps the list at 10 lines.
\subject{Talks}
% No back-to-back roadmaps (Yiqing, 2026-10-11): when a subsection starts right
% after its section's roadmap, with no content slide in between, its own
% roadmap would repeat the one just shown (5-1's handout showed "Two Cases"
% twice), so it is skipped. The section roadmap records the frame count; the
% subsection roadmap appears only if a content slide has come since.
\makeatletter
\def\course@secframe{-1}
\AtBeginSection[]
{
  \begin{frame}[plain,c]
    \frametitle{}
    \tableofcontents[currentsection,hideothersubsections]
    \addtocounter{framenumber}{-1}
  \end{frame}
  \xdef\course@secframe{\the\c@framenumber}
}
\AtBeginSubsection[]
{
  \ifnum\c@framenumber=\course@secframe\relax\else
  \begin{frame}[plain,c]
    \frametitle{}
    \tableofcontents[currentsection,currentsubsection,hideothersubsections]
    \addtocounter{framenumber}{-1}
  \end{frame}
  \fi
}
\makeatother

% ------------------------------------------------------ theorem environments
\newtheorem{assumption}{Assumption}
\newtheorem{iass}{Identification Assumption}
\newtheorem{ires}{Identification Result}
\newtheorem{estm}{Estimand}
\newtheorem{esti}{Estimator}
\newtheorem{question}{Question}
\newtheorem{answer}{Answer}
\newtheorem{proposition}{Proposition}

% =====================================================================
% Notation
% =====================================================================
\definecolor{dgreen}{rgb}{0,.6,0.8}

\renewcommand\r{\right}
\renewcommand\l{\left}
\newcommand{\bs}{\boldsymbol}
\newcommand{\mb}{\mathbf}
\newcommand{\real}{\mathbb{R}}
\newcommand{\cN}{\mathcal{N}}
\newcommand{\E}{\mathbb{E}}
\newcommand{\V}{\mathbb{V}}

% MERGE NOTE 1. Lecture 2-4 set \epsilon to \varepsilon; the other twelve
% left it at the upright default. Adopting \varepsilon course-wide so the
% error term looks the same in every lecture.
\renewcommand{\epsilon}{\varepsilon}

% MERGE NOTE 2. Lecture 2-4 wrote its vectors as \bm (bold italic); the other
% twelve use \mathbf (bold upright). \mathbf wins --- it is the convention in
% eleven of the thirteen decks and in the course's matrix lectures. Only
% \iid, \E, \cN and \Var from that block are actually used in 2-4's body.
\renewcommand{\u}{\ensuremath{\mb{u}}}
\renewcommand{\v}{\ensuremath{\mb{v}}}
\newcommand{\U}{\ensuremath{\mb{U}}}
\newcommand{\M}{\ensuremath{\mb{M}}}
\newcommand{\X}{\ensuremath{\mb{X}}}
\newcommand{\x}{\ensuremath{\mb{x}}}
\newcommand{\y}{\ensuremath{\mb{y}}}
\newcommand{\Z}{\ensuremath{\mb{Z}}}
\newcommand{\C}{\ensuremath{\mb{C}}}
\newcommand{\D}{\ensuremath{\mb{D}}}
\newcommand{\I}{\ensuremath{\mb{I}}}
\newcommand{\W}{\ensuremath{\mb{W}}}
\renewcommand{\a}{\ensuremath{\mb{a}}}
\renewcommand{\P}{\ensuremath{\mb{P}}}
\newcommand{\one}{\ensuremath{\mb{1}}}
\newcommand{\zero}{\ensuremath{\mb{0}}}

\renewcommand{\b}{\ensuremath{\bs{\beta}}}
\newcommand{\be}{\ensuremath{\bs{\beta}}}
\newcommand{\e}{\ensuremath{\bs{\epsilon}}}
\newcommand{\eps}{\ensuremath{\bs{\epsilon}}}
\newcommand{\eeta}{\ensuremath{\bs{\eta}}}
\newcommand{\bhat}{\ensuremath{\hat{\bs{\beta}}}}
\newcommand{\ehat}{\ensuremath{\hat{\e}}}
\newcommand{\Sig}{\ensuremath{\bs{\Sigma}}}
\newcommand{\SigInv}{\ensuremath{\bs{\Sigma^{-1}}}}

\newcommand{\XX}{\ensuremath{\X'\X}}
\newcommand{\XXinv}{\ensuremath{\left(\XX\right)^{-1}}}
\newcommand{\hatmat}{\ensuremath{\X \XXinv \X'}}
\newcommand{\hatsig}{\ensuremath{\hat{\sigma}^2}}
\newcommand{\sig}{\ensuremath{\sigma^2}}

\newcommand{\cA}{\mathcal{A}}
\newcommand{\cJ}{\mathcal{J}}
\newcommand{\cL}{\mathcal{L}}
\newcommand{\cLL}{\mathcal{LL}}
\newcommand{\cll}{\ell}
\newcommand{\cQ}{\mathcal{Q}}
\newcommand{\cP}{\mathcal{P}}

% MERGE NOTE 3. \var / \Var and \cov / \Cov were spelled with \text in some
% lectures and \mbox in others. Same output; \text kept.
\newcommand{\var}{\text{Var}}
\newcommand{\Var}{\text{Var}}
\newcommand{\cov}{\text{Cov}}
\newcommand{\Cov}{\text{Cov}}
\newcommand{\cor}{\text{Cor}}
\newcommand{\tr}{\text{tr}}
\newcommand{\avepop}{\frac{1}{n}\sum_{i=1}^n}

\DeclareMathOperator*\plim{plim}
\DeclareMathOperator\rank{rank}
\DeclareMathOperator*\argmax{argmax}
\DeclareMathOperator*\argmin{argmin}
\DeclareMathOperator{\sgn}{sgn}
\newcommand{\diag}{\mathop{\mathrm{diag}}}

\newcommand\ud{\mathrm{d}}
\newcommand\dist{\buildrel\rm d\over\sim}
\newcommand\ind{\stackrel{\rm indep.}{\sim}}
\newcommand\iid{\stackrel{\rm i.i.d.}{\sim}}
\newcommand\logit{{\rm logit}}

\newcommand\indep{\protect\mathpalette{\protect\independenT}{\perp}}
\def\independenT#1#2{\mathrel{\rlap{$#1#2$}\mkern2mu{#1#2}}}

% ------------------------------------------------------------- text helpers
\newcommand{\bi}{\begin{itemize}}
\newcommand{\ei}{\end{itemize}}
\newcommand{\noi}{\noindent}
\newcommand{\red}{\textcolor{Cardinal}}
\newcommand{\blue}{\textcolor{CourseBlue}}
\newcommand{\gray}{\textcolor{gray}}
\newcommand\spacingset[1]{\renewcommand{\baselinestretch}{#1}\small\normalsize}

\hyphenation{tyamamot}


% ------------------------------------------------ figure-position check (2026-10-11)
% Records where every \includegraphics lands and its size (a whatsit; it moves
% nothing). check_slides.py flags two consecutive slides that show a figure of
% the same size in different places: the class sees it jump (Lecture 3, slides
% 18 to 19: one frame had \vspace{.3cm} above the figure, the next did not),
% and two consecutive slides that show the same file at different sizes or
% places (Lecture 6-3, slides 10 to 11, the ballot image). The file name is the
% seventh field.
\makeatletter
\providecommand\figpos[7]{}
\AtBeginDocument{\immediate\write\@auxout{\string\providecommand\string\figpos[7]{}}}
\newbox\course@figbox
\let\course@origGinclude\Ginclude@graphics
\def\Ginclude@graphics#1{%
  \setbox\course@figbox\hbox{\course@origGinclude{#1}}%
  \pdfsavepos
  \edef\course@figtmp{\noexpand\write\noexpand\@auxout{\string\figpos
    {\noexpand\thepage}{\noexpand\insertframenumber}%
    {\noexpand\the\noexpand\pdflastxpos}{\noexpand\the\noexpand\pdflastypos}%
    {\number\wd\course@figbox}{\number\ht\course@figbox}{\detokenize{#1}}}}%
  \course@figtmp
  \box\course@figbox}
\makeatother

%   % =====================================================================
% POLI 450A (2026) --- shared visual system.
% Input from each lecture AFTER ../course_preamble.tex.
%
% Same shape as the 2026 Northwestern workshop theme
% (Teaching/PanelData/_latex/2026/202608_Northwestern/workshop_theme.tex),
% with the workshop's purple swapped for the course's cardinal. Kept
% deliberately light so the lecture sources stay easy to edit.
%
% To retune the course to Stanford's official cardinal, change the one
% Cardinal line below to HTML 8C1515.
% =====================================================================
\definecolor{Cardinal}{HTML}{A8003B}     % primary: titles, structure, alerts
\definecolor{CourseInk}{HTML}{252A34}    % body text
\definecolor{CourseBlue}{HTML}{2F6B9A}   % secondary accent
\definecolor{CourseMist}{HTML}{F8F1F4}   % block backgrounds

\setbeamercolor{normal text}{fg=CourseInk,bg=white}
\setbeamercolor{structure}{fg=Cardinal}
\setbeamercolor{alerted text}{fg=Cardinal}
\setbeamercolor{frametitle}{fg=white,bg=Cardinal}
\setbeamercolor{title}{fg=Cardinal}
\setbeamercolor{subtitle}{fg=CourseInk}
\setbeamercolor{block title}{fg=white,bg=Cardinal}
\setbeamercolor{block body}{fg=CourseInk,bg=CourseMist}
\setbeamercolor{section in toc}{fg=Cardinal}
\setbeamercolor{subsection in toc}{fg=CourseInk}
\setbeamercolor{itemize item}{fg=Cardinal}
\setbeamercolor{itemize subitem}{fg=Cardinal}
\setbeamercolor{enumerate item}{fg=Cardinal}

\setbeamerfont{title}{series=\bfseries,size=\Large}
\setbeamerfont{frametitle}{series=\bfseries,size=\large}
\setbeamerfont{block title}{series=\bfseries}
\setbeamertemplate{navigation symbols}{}
\setbeamertemplate{items}[circle]
\setbeamertemplate{blocks}[rounded][shadow=false]
\setbeamertemplate{frametitle continuation}{}
\setbeamertemplate{footline}{%
  \leavevmode\hbox{%
  \begin{beamercolorbox}[wd=.88\paperwidth,ht=2.6ex,dp=1.2ex,leftskip=1.2em]{author in head/foot}%
  \usebeamerfont{author in head/foot}\insertshorttitle
  \end{beamercolorbox}%
  \begin{beamercolorbox}[wd=.12\paperwidth,ht=2.6ex,dp=1.2ex,center]{date in head/foot}%
  \insertframenumber
  \end{beamercolorbox}}}

\newcommand{\slidesource}[1]{\vfill{\scriptsize\color{gray}Source: #1}}
\newcommand{\keyquestion}[1]{\begin{block}{Question}#1\end{block}}

%   \graphicspath{{./}{./figs/}}
%   \title[Short title]{450A Political Methodology I}
%   \subtitle{Lecture N: ...}
%
% Modeled on the 2026 Northwestern workshop template
% (Teaching/PanelData/_latex/2026/202608_Northwestern/workshop_preamble.tex).
%
% SCOPE: theme, beamer layout, packages, the roadmap hooks, and --- unlike
% the Northwestern decks --- the notation macros. The 450A lectures
% grew from one common ancestor, so their macro blocks agree wherever they
% overlap; merging them is safe. Three points where they had drifted are
% resolved here, and each is flagged inline below.
% =====================================================================

\usetheme{metropolis}
\metroset{sectionpage=none, subsectionpage=none, numbering=none}
\setbeamertemplate{navigation symbols}{}

\setbeamersize{text margin left=5mm, text margin right=5mm}
\setbeamertemplate{itemize/enumerate body begin}{\leftmargini=1.5em\relax}
\setbeamertemplate{itemize/enumerate subbody begin}{\leftmarginii=1.5em\relax}
\setbeamertemplate{itemize/enumerate subsubbody begin}{\leftmarginiii=1.5em\relax}

% ----------------------------------------------------------- top alignment
% Every deck passes the t class option, so beamer applies the frame key
% below to every frame (explicit [plain,c] frames — title page, roadmaps —
% stay centered). Beamer's stock t key puts a rigid .2cm pad above the
% content; frames sized to this template's full-height ceilings then run
% Overfull \vbox by up to that pad. Same key as beamerbaseframe.sty, with
% one change: the pad may shrink to zero when the frame is full.
\makeatletter
\define@key{beamerframe}{t}[true]{% top
  \beamer@frametopskip=.2cm plus .5\paperheight minus .2cm\relax%
  \beamer@framebottomskip=0pt plus 1fill\relax%
  \beamer@frametopskipautobreak=\beamer@frametopskip\relax%
  \beamer@framebottomskipautobreak=\beamer@framebottomskip\relax%
}

% ------------------------------------------------- slide-top check (2026-10-11)
% A slide whose content is too tall for the frame is squeezed: the pad above
% shrinks and its content starts higher than on other slides, so a figure
% "wobbles" when the class steps from one slide to the next (Lecture 4-1,
% slides 1-2, 2026-10-11). The marker below records where each slide's
% content starts (a whatsit: it adds no space and moves nothing).
% `python3 ../check_slides.py lec<N>` then lists every slide that starts
% higher than the deck's normal position. That list must be empty.
% Recorded per slide: page, frame number, content-top y (sp), the room left
% below the title bar, and P for a plain frame (title page, roadmaps).
\providecommand\slidetop[5]{}
\AtBeginDocument{\immediate\write\@auxout{\string\providecommand\string\slidetop[5]{}}}
\def\beamer@markframetop{%
  \pdfsavepos
  \edef\beamer@marktmp{\noexpand\write\noexpand\@auxout{\string\slidetop
    {\noexpand\thepage}{\noexpand\insertframenumber}{\noexpand\the\noexpand\pdflastypos}%
    {\number\beamer@frametextheight}{\ifbeamer@plainframe P\else N\fi}}}%
  \beamer@marktmp}
\RequirePackage{etoolbox}
\patchcmd{\endbeamer@frameslide}
  {\vskip\beamer@frametopskip}
  {\vskip\beamer@frametopskip\beamer@markframetop}
  {}{\PackageWarning{course_preamble}{slide-top marker not installed}}
\makeatother

% ---------------------------------------------------------------- packages
\usepackage{xcolor}
\usepackage{colortbl}
\usepackage{amsmath}
\usepackage{amsfonts}
\usepackage{amssymb}
\usepackage{bm}
\usepackage{graphicx}
\usepackage{multirow}
\usepackage{fancyvrb}
\usepackage{listings}
\usepackage{tikz}
\usetikzlibrary{arrows.meta,positioning,decorations.pathreplacing}

% --------------------------------------------------------- figure fitting
% Default bounding box for every \includegraphics that does not override it.
% With keepaspectratio the image is scaled to fit *inside* width x height, so
% a figure can never run past the side margins or off the bottom of a slide.
% Options given at the call site override these on a per-key basis: e.g.
% [width=.6\linewidth] narrows the box but keeps the height ceiling.
% \figmaxheight is set at \begin{document}: beamer only computes the final
% \textheight once the head/foot templates are known. 0.80 is just under what a
% frame with a title and one line of text can hold (measured: 0.834\textheight);
% a frame that needs more room for text can lower it locally.
\newlength{\figmaxheight}
\AtBeginDocument{\setlength{\figmaxheight}{0.80\textheight}}
\setkeys{Gin}{width=\linewidth,height=\figmaxheight,keepaspectratio}

% -------------------------------------------------------- roadmap frames
% Identical in all lectures. A lecture that does not want the
% subsection roadmap clears it after the \input with \AtBeginSubsection[]{}.
% hideothersubsections is load-bearing, not cosmetic: lecture 4 had 13
% subsections before it was split into 4-1, 4-2 and 4-3 (2026-10-09), and a
% full contents list ran off the top and bottom of the slide. Hiding the other sections' subsections caps the list at 10 lines.
\subject{Talks}
% No back-to-back roadmaps (Yiqing, 2026-10-11): when a subsection starts right
% after its section's roadmap, with no content slide in between, its own
% roadmap would repeat the one just shown (5-1's handout showed "Two Cases"
% twice), so it is skipped. The section roadmap records the frame count; the
% subsection roadmap appears only if a content slide has come since.
\makeatletter
\def\course@secframe{-1}
\AtBeginSection[]
{
  \begin{frame}[plain,c]
    \frametitle{}
    \tableofcontents[currentsection,hideothersubsections]
    \addtocounter{framenumber}{-1}
  \end{frame}
  \xdef\course@secframe{\the\c@framenumber}
}
\AtBeginSubsection[]
{
  \ifnum\c@framenumber=\course@secframe\relax\else
  \begin{frame}[plain,c]
    \frametitle{}
    \tableofcontents[currentsection,currentsubsection,hideothersubsections]
    \addtocounter{framenumber}{-1}
  \end{frame}
  \fi
}
\makeatother

% ------------------------------------------------------ theorem environments
\newtheorem{assumption}{Assumption}
\newtheorem{iass}{Identification Assumption}
\newtheorem{ires}{Identification Result}
\newtheorem{estm}{Estimand}
\newtheorem{esti}{Estimator}
\newtheorem{question}{Question}
\newtheorem{answer}{Answer}
\newtheorem{proposition}{Proposition}

% =====================================================================
% Notation
% =====================================================================
\definecolor{dgreen}{rgb}{0,.6,0.8}

\renewcommand\r{\right}
\renewcommand\l{\left}
\newcommand{\bs}{\boldsymbol}
\newcommand{\mb}{\mathbf}
\newcommand{\real}{\mathbb{R}}
\newcommand{\cN}{\mathcal{N}}
\newcommand{\E}{\mathbb{E}}
\newcommand{\V}{\mathbb{V}}

% MERGE NOTE 1. Lecture 2-4 set \epsilon to \varepsilon; the other twelve
% left it at the upright default. Adopting \varepsilon course-wide so the
% error term looks the same in every lecture.
\renewcommand{\epsilon}{\varepsilon}

% MERGE NOTE 2. Lecture 2-4 wrote its vectors as \bm (bold italic); the other
% twelve use \mathbf (bold upright). \mathbf wins --- it is the convention in
% eleven of the thirteen decks and in the course's matrix lectures. Only
% \iid, \E, \cN and \Var from that block are actually used in 2-4's body.
\renewcommand{\u}{\ensuremath{\mb{u}}}
\renewcommand{\v}{\ensuremath{\mb{v}}}
\newcommand{\U}{\ensuremath{\mb{U}}}
\newcommand{\M}{\ensuremath{\mb{M}}}
\newcommand{\X}{\ensuremath{\mb{X}}}
\newcommand{\x}{\ensuremath{\mb{x}}}
\newcommand{\y}{\ensuremath{\mb{y}}}
\newcommand{\Z}{\ensuremath{\mb{Z}}}
\newcommand{\C}{\ensuremath{\mb{C}}}
\newcommand{\D}{\ensuremath{\mb{D}}}
\newcommand{\I}{\ensuremath{\mb{I}}}
\newcommand{\W}{\ensuremath{\mb{W}}}
\renewcommand{\a}{\ensuremath{\mb{a}}}
\renewcommand{\P}{\ensuremath{\mb{P}}}
\newcommand{\one}{\ensuremath{\mb{1}}}
\newcommand{\zero}{\ensuremath{\mb{0}}}

\renewcommand{\b}{\ensuremath{\bs{\beta}}}
\newcommand{\be}{\ensuremath{\bs{\beta}}}
\newcommand{\e}{\ensuremath{\bs{\epsilon}}}
\newcommand{\eps}{\ensuremath{\bs{\epsilon}}}
\newcommand{\eeta}{\ensuremath{\bs{\eta}}}
\newcommand{\bhat}{\ensuremath{\hat{\bs{\beta}}}}
\newcommand{\ehat}{\ensuremath{\hat{\e}}}
\newcommand{\Sig}{\ensuremath{\bs{\Sigma}}}
\newcommand{\SigInv}{\ensuremath{\bs{\Sigma^{-1}}}}

\newcommand{\XX}{\ensuremath{\X'\X}}
\newcommand{\XXinv}{\ensuremath{\left(\XX\right)^{-1}}}
\newcommand{\hatmat}{\ensuremath{\X \XXinv \X'}}
\newcommand{\hatsig}{\ensuremath{\hat{\sigma}^2}}
\newcommand{\sig}{\ensuremath{\sigma^2}}

\newcommand{\cA}{\mathcal{A}}
\newcommand{\cJ}{\mathcal{J}}
\newcommand{\cL}{\mathcal{L}}
\newcommand{\cLL}{\mathcal{LL}}
\newcommand{\cll}{\ell}
\newcommand{\cQ}{\mathcal{Q}}
\newcommand{\cP}{\mathcal{P}}

% MERGE NOTE 3. \var / \Var and \cov / \Cov were spelled with \text in some
% lectures and \mbox in others. Same output; \text kept.
\newcommand{\var}{\text{Var}}
\newcommand{\Var}{\text{Var}}
\newcommand{\cov}{\text{Cov}}
\newcommand{\Cov}{\text{Cov}}
\newcommand{\cor}{\text{Cor}}
\newcommand{\tr}{\text{tr}}
\newcommand{\avepop}{\frac{1}{n}\sum_{i=1}^n}

\DeclareMathOperator*\plim{plim}
\DeclareMathOperator\rank{rank}
\DeclareMathOperator*\argmax{argmax}
\DeclareMathOperator*\argmin{argmin}
\DeclareMathOperator{\sgn}{sgn}
\newcommand{\diag}{\mathop{\mathrm{diag}}}

\newcommand\ud{\mathrm{d}}
\newcommand\dist{\buildrel\rm d\over\sim}
\newcommand\ind{\stackrel{\rm indep.}{\sim}}
\newcommand\iid{\stackrel{\rm i.i.d.}{\sim}}
\newcommand\logit{{\rm logit}}

\newcommand\indep{\protect\mathpalette{\protect\independenT}{\perp}}
\def\independenT#1#2{\mathrel{\rlap{$#1#2$}\mkern2mu{#1#2}}}

% ------------------------------------------------------------- text helpers
\newcommand{\bi}{\begin{itemize}}
\newcommand{\ei}{\end{itemize}}
\newcommand{\noi}{\noindent}
\newcommand{\red}{\textcolor{Cardinal}}
\newcommand{\blue}{\textcolor{CourseBlue}}
\newcommand{\gray}{\textcolor{gray}}
\newcommand\spacingset[1]{\renewcommand{\baselinestretch}{#1}\small\normalsize}

\hyphenation{tyamamot}


% ------------------------------------------------ figure-position check (2026-10-11)
% Records where every \includegraphics lands and its size (a whatsit; it moves
% nothing). check_slides.py flags two consecutive slides that show a figure of
% the same size in different places: the class sees it jump (Lecture 3, slides
% 18 to 19: one frame had \vspace{.3cm} above the figure, the next did not),
% and two consecutive slides that show the same file at different sizes or
% places (Lecture 6-3, slides 10 to 11, the ballot image). The file name is the
% seventh field.
\makeatletter
\providecommand\figpos[7]{}
\AtBeginDocument{\immediate\write\@auxout{\string\providecommand\string\figpos[7]{}}}
\newbox\course@figbox
\let\course@origGinclude\Ginclude@graphics
\def\Ginclude@graphics#1{%
  \setbox\course@figbox\hbox{\course@origGinclude{#1}}%
  \pdfsavepos
  \edef\course@figtmp{\noexpand\write\noexpand\@auxout{\string\figpos
    {\noexpand\thepage}{\noexpand\insertframenumber}%
    {\noexpand\the\noexpand\pdflastxpos}{\noexpand\the\noexpand\pdflastypos}%
    {\number\wd\course@figbox}{\number\ht\course@figbox}{\detokenize{#1}}}}%
  \course@figtmp
  \box\course@figbox}
\makeatother

%   % =====================================================================
% POLI 450A (2026) --- shared visual system.
% Input from each lecture AFTER ../course_preamble.tex.
%
% Same shape as the 2026 Northwestern workshop theme
% (Teaching/PanelData/_latex/2026/202608_Northwestern/workshop_theme.tex),
% with the workshop's purple swapped for the course's cardinal. Kept
% deliberately light so the lecture sources stay easy to edit.
%
% To retune the course to Stanford's official cardinal, change the one
% Cardinal line below to HTML 8C1515.
% =====================================================================
\definecolor{Cardinal}{HTML}{A8003B}     % primary: titles, structure, alerts
\definecolor{CourseInk}{HTML}{252A34}    % body text
\definecolor{CourseBlue}{HTML}{2F6B9A}   % secondary accent
\definecolor{CourseMist}{HTML}{F8F1F4}   % block backgrounds

\setbeamercolor{normal text}{fg=CourseInk,bg=white}
\setbeamercolor{structure}{fg=Cardinal}
\setbeamercolor{alerted text}{fg=Cardinal}
\setbeamercolor{frametitle}{fg=white,bg=Cardinal}
\setbeamercolor{title}{fg=Cardinal}
\setbeamercolor{subtitle}{fg=CourseInk}
\setbeamercolor{block title}{fg=white,bg=Cardinal}
\setbeamercolor{block body}{fg=CourseInk,bg=CourseMist}
\setbeamercolor{section in toc}{fg=Cardinal}
\setbeamercolor{subsection in toc}{fg=CourseInk}
\setbeamercolor{itemize item}{fg=Cardinal}
\setbeamercolor{itemize subitem}{fg=Cardinal}
\setbeamercolor{enumerate item}{fg=Cardinal}

\setbeamerfont{title}{series=\bfseries,size=\Large}
\setbeamerfont{frametitle}{series=\bfseries,size=\large}
\setbeamerfont{block title}{series=\bfseries}
\setbeamertemplate{navigation symbols}{}
\setbeamertemplate{items}[circle]
\setbeamertemplate{blocks}[rounded][shadow=false]
\setbeamertemplate{frametitle continuation}{}
\setbeamertemplate{footline}{%
  \leavevmode\hbox{%
  \begin{beamercolorbox}[wd=.88\paperwidth,ht=2.6ex,dp=1.2ex,leftskip=1.2em]{author in head/foot}%
  \usebeamerfont{author in head/foot}\insertshorttitle
  \end{beamercolorbox}%
  \begin{beamercolorbox}[wd=.12\paperwidth,ht=2.6ex,dp=1.2ex,center]{date in head/foot}%
  \insertframenumber
  \end{beamercolorbox}}}

\newcommand{\slidesource}[1]{\vfill{\scriptsize\color{gray}Source: #1}}
\newcommand{\keyquestion}[1]{\begin{block}{Question}#1\end{block}}

%   \graphicspath{{./}{./figs/}}
%   \title[Short title]{450A Political Methodology I}
%   \subtitle{Lecture N: ...}
%
% Modeled on the 2026 Northwestern workshop template
% (Teaching/PanelData/_latex/2026/202608_Northwestern/workshop_preamble.tex).
%
% SCOPE: theme, beamer layout, packages, the roadmap hooks, and --- unlike
% the Northwestern decks --- the notation macros. The 450A lectures
% grew from one common ancestor, so their macro blocks agree wherever they
% overlap; merging them is safe. Three points where they had drifted are
% resolved here, and each is flagged inline below.
% =====================================================================

\usetheme{metropolis}
\metroset{sectionpage=none, subsectionpage=none, numbering=none}
\setbeamertemplate{navigation symbols}{}

\setbeamersize{text margin left=5mm, text margin right=5mm}
\setbeamertemplate{itemize/enumerate body begin}{\leftmargini=1.5em\relax}
\setbeamertemplate{itemize/enumerate subbody begin}{\leftmarginii=1.5em\relax}
\setbeamertemplate{itemize/enumerate subsubbody begin}{\leftmarginiii=1.5em\relax}

% ----------------------------------------------------------- top alignment
% Every deck passes the t class option, so beamer applies the frame key
% below to every frame (explicit [plain,c] frames — title page, roadmaps —
% stay centered). Beamer's stock t key puts a rigid .2cm pad above the
% content; frames sized to this template's full-height ceilings then run
% Overfull \vbox by up to that pad. Same key as beamerbaseframe.sty, with
% one change: the pad may shrink to zero when the frame is full.
\makeatletter
\define@key{beamerframe}{t}[true]{% top
  \beamer@frametopskip=.2cm plus .5\paperheight minus .2cm\relax%
  \beamer@framebottomskip=0pt plus 1fill\relax%
  \beamer@frametopskipautobreak=\beamer@frametopskip\relax%
  \beamer@framebottomskipautobreak=\beamer@framebottomskip\relax%
}

% ------------------------------------------------- slide-top check (2026-10-11)
% A slide whose content is too tall for the frame is squeezed: the pad above
% shrinks and its content starts higher than on other slides, so a figure
% "wobbles" when the class steps from one slide to the next (Lecture 4-1,
% slides 1-2, 2026-10-11). The marker below records where each slide's
% content starts (a whatsit: it adds no space and moves nothing).
% `python3 ../check_slides.py lec<N>` then lists every slide that starts
% higher than the deck's normal position. That list must be empty.
% Recorded per slide: page, frame number, content-top y (sp), the room left
% below the title bar, and P for a plain frame (title page, roadmaps).
\providecommand\slidetop[5]{}
\AtBeginDocument{\immediate\write\@auxout{\string\providecommand\string\slidetop[5]{}}}
\def\beamer@markframetop{%
  \pdfsavepos
  \edef\beamer@marktmp{\noexpand\write\noexpand\@auxout{\string\slidetop
    {\noexpand\thepage}{\noexpand\insertframenumber}{\noexpand\the\noexpand\pdflastypos}%
    {\number\beamer@frametextheight}{\ifbeamer@plainframe P\else N\fi}}}%
  \beamer@marktmp}
\RequirePackage{etoolbox}
\patchcmd{\endbeamer@frameslide}
  {\vskip\beamer@frametopskip}
  {\vskip\beamer@frametopskip\beamer@markframetop}
  {}{\PackageWarning{course_preamble}{slide-top marker not installed}}
\makeatother

% ---------------------------------------------------------------- packages
\usepackage{xcolor}
\usepackage{colortbl}
\usepackage{amsmath}
\usepackage{amsfonts}
\usepackage{amssymb}
\usepackage{bm}
\usepackage{graphicx}
\usepackage{multirow}
\usepackage{fancyvrb}
\usepackage{listings}
\usepackage{tikz}
\usetikzlibrary{arrows.meta,positioning,decorations.pathreplacing}

% --------------------------------------------------------- figure fitting
% Default bounding box for every \includegraphics that does not override it.
% With keepaspectratio the image is scaled to fit *inside* width x height, so
% a figure can never run past the side margins or off the bottom of a slide.
% Options given at the call site override these on a per-key basis: e.g.
% [width=.6\linewidth] narrows the box but keeps the height ceiling.
% \figmaxheight is set at \begin{document}: beamer only computes the final
% \textheight once the head/foot templates are known. 0.80 is just under what a
% frame with a title and one line of text can hold (measured: 0.834\textheight);
% a frame that needs more room for text can lower it locally.
\newlength{\figmaxheight}
\AtBeginDocument{\setlength{\figmaxheight}{0.80\textheight}}
\setkeys{Gin}{width=\linewidth,height=\figmaxheight,keepaspectratio}

% -------------------------------------------------------- roadmap frames
% Identical in all lectures. A lecture that does not want the
% subsection roadmap clears it after the \input with \AtBeginSubsection[]{}.
% hideothersubsections is load-bearing, not cosmetic: lecture 4 had 13
% subsections before it was split into 4-1, 4-2 and 4-3 (2026-10-09), and a
% full contents list ran off the top and bottom of the slide. Hiding the other sections' subsections caps the list at 10 lines.
\subject{Talks}
% No back-to-back roadmaps (Yiqing, 2026-10-11): when a subsection starts right
% after its section's roadmap, with no content slide in between, its own
% roadmap would repeat the one just shown (5-1's handout showed "Two Cases"
% twice), so it is skipped. The section roadmap records the frame count; the
% subsection roadmap appears only if a content slide has come since.
\makeatletter
\def\course@secframe{-1}
\AtBeginSection[]
{
  \begin{frame}[plain,c]
    \frametitle{}
    \tableofcontents[currentsection,hideothersubsections]
    \addtocounter{framenumber}{-1}
  \end{frame}
  \xdef\course@secframe{\the\c@framenumber}
}
\AtBeginSubsection[]
{
  \ifnum\c@framenumber=\course@secframe\relax\else
  \begin{frame}[plain,c]
    \frametitle{}
    \tableofcontents[currentsection,currentsubsection,hideothersubsections]
    \addtocounter{framenumber}{-1}
  \end{frame}
  \fi
}
\makeatother

% ------------------------------------------------------ theorem environments
\newtheorem{assumption}{Assumption}
\newtheorem{iass}{Identification Assumption}
\newtheorem{ires}{Identification Result}
\newtheorem{estm}{Estimand}
\newtheorem{esti}{Estimator}
\newtheorem{question}{Question}
\newtheorem{answer}{Answer}
\newtheorem{proposition}{Proposition}

% =====================================================================
% Notation
% =====================================================================
\definecolor{dgreen}{rgb}{0,.6,0.8}

\renewcommand\r{\right}
\renewcommand\l{\left}
\newcommand{\bs}{\boldsymbol}
\newcommand{\mb}{\mathbf}
\newcommand{\real}{\mathbb{R}}
\newcommand{\cN}{\mathcal{N}}
\newcommand{\E}{\mathbb{E}}
\newcommand{\V}{\mathbb{V}}

% MERGE NOTE 1. Lecture 2-4 set \epsilon to \varepsilon; the other twelve
% left it at the upright default. Adopting \varepsilon course-wide so the
% error term looks the same in every lecture.
\renewcommand{\epsilon}{\varepsilon}

% MERGE NOTE 2. Lecture 2-4 wrote its vectors as \bm (bold italic); the other
% twelve use \mathbf (bold upright). \mathbf wins --- it is the convention in
% eleven of the thirteen decks and in the course's matrix lectures. Only
% \iid, \E, \cN and \Var from that block are actually used in 2-4's body.
\renewcommand{\u}{\ensuremath{\mb{u}}}
\renewcommand{\v}{\ensuremath{\mb{v}}}
\newcommand{\U}{\ensuremath{\mb{U}}}
\newcommand{\M}{\ensuremath{\mb{M}}}
\newcommand{\X}{\ensuremath{\mb{X}}}
\newcommand{\x}{\ensuremath{\mb{x}}}
\newcommand{\y}{\ensuremath{\mb{y}}}
\newcommand{\Z}{\ensuremath{\mb{Z}}}
\newcommand{\C}{\ensuremath{\mb{C}}}
\newcommand{\D}{\ensuremath{\mb{D}}}
\newcommand{\I}{\ensuremath{\mb{I}}}
\newcommand{\W}{\ensuremath{\mb{W}}}
\renewcommand{\a}{\ensuremath{\mb{a}}}
\renewcommand{\P}{\ensuremath{\mb{P}}}
\newcommand{\one}{\ensuremath{\mb{1}}}
\newcommand{\zero}{\ensuremath{\mb{0}}}

\renewcommand{\b}{\ensuremath{\bs{\beta}}}
\newcommand{\be}{\ensuremath{\bs{\beta}}}
\newcommand{\e}{\ensuremath{\bs{\epsilon}}}
\newcommand{\eps}{\ensuremath{\bs{\epsilon}}}
\newcommand{\eeta}{\ensuremath{\bs{\eta}}}
\newcommand{\bhat}{\ensuremath{\hat{\bs{\beta}}}}
\newcommand{\ehat}{\ensuremath{\hat{\e}}}
\newcommand{\Sig}{\ensuremath{\bs{\Sigma}}}
\newcommand{\SigInv}{\ensuremath{\bs{\Sigma^{-1}}}}

\newcommand{\XX}{\ensuremath{\X'\X}}
\newcommand{\XXinv}{\ensuremath{\left(\XX\right)^{-1}}}
\newcommand{\hatmat}{\ensuremath{\X \XXinv \X'}}
\newcommand{\hatsig}{\ensuremath{\hat{\sigma}^2}}
\newcommand{\sig}{\ensuremath{\sigma^2}}

\newcommand{\cA}{\mathcal{A}}
\newcommand{\cJ}{\mathcal{J}}
\newcommand{\cL}{\mathcal{L}}
\newcommand{\cLL}{\mathcal{LL}}
\newcommand{\cll}{\ell}
\newcommand{\cQ}{\mathcal{Q}}
\newcommand{\cP}{\mathcal{P}}

% MERGE NOTE 3. \var / \Var and \cov / \Cov were spelled with \text in some
% lectures and \mbox in others. Same output; \text kept.
\newcommand{\var}{\text{Var}}
\newcommand{\Var}{\text{Var}}
\newcommand{\cov}{\text{Cov}}
\newcommand{\Cov}{\text{Cov}}
\newcommand{\cor}{\text{Cor}}
\newcommand{\tr}{\text{tr}}
\newcommand{\avepop}{\frac{1}{n}\sum_{i=1}^n}

\DeclareMathOperator*\plim{plim}
\DeclareMathOperator\rank{rank}
\DeclareMathOperator*\argmax{argmax}
\DeclareMathOperator*\argmin{argmin}
\DeclareMathOperator{\sgn}{sgn}
\newcommand{\diag}{\mathop{\mathrm{diag}}}

\newcommand\ud{\mathrm{d}}
\newcommand\dist{\buildrel\rm d\over\sim}
\newcommand\ind{\stackrel{\rm indep.}{\sim}}
\newcommand\iid{\stackrel{\rm i.i.d.}{\sim}}
\newcommand\logit{{\rm logit}}

\newcommand\indep{\protect\mathpalette{\protect\independenT}{\perp}}
\def\independenT#1#2{\mathrel{\rlap{$#1#2$}\mkern2mu{#1#2}}}

% ------------------------------------------------------------- text helpers
\newcommand{\bi}{\begin{itemize}}
\newcommand{\ei}{\end{itemize}}
\newcommand{\noi}{\noindent}
\newcommand{\red}{\textcolor{Cardinal}}
\newcommand{\blue}{\textcolor{CourseBlue}}
\newcommand{\gray}{\textcolor{gray}}
\newcommand\spacingset[1]{\renewcommand{\baselinestretch}{#1}\small\normalsize}

\hyphenation{tyamamot}


% ------------------------------------------------ figure-position check (2026-10-11)
% Records where every \includegraphics lands and its size (a whatsit; it moves
% nothing). check_slides.py flags two consecutive slides that show a figure of
% the same size in different places: the class sees it jump (Lecture 3, slides
% 18 to 19: one frame had \vspace{.3cm} above the figure, the next did not),
% and two consecutive slides that show the same file at different sizes or
% places (Lecture 6-3, slides 10 to 11, the ballot image). The file name is the
% seventh field.
\makeatletter
\providecommand\figpos[7]{}
\AtBeginDocument{\immediate\write\@auxout{\string\providecommand\string\figpos[7]{}}}
\newbox\course@figbox
\let\course@origGinclude\Ginclude@graphics
\def\Ginclude@graphics#1{%
  \setbox\course@figbox\hbox{\course@origGinclude{#1}}%
  \pdfsavepos
  \edef\course@figtmp{\noexpand\write\noexpand\@auxout{\string\figpos
    {\noexpand\thepage}{\noexpand\insertframenumber}%
    {\noexpand\the\noexpand\pdflastxpos}{\noexpand\the\noexpand\pdflastypos}%
    {\number\wd\course@figbox}{\number\ht\course@figbox}{\detokenize{#1}}}}%
  \course@figtmp
  \box\course@figbox}
\makeatother
   % theme, layout, packages, notation
% =====================================================================
% POLI 450A (2026) --- shared visual system.
% Input from each lecture AFTER ../course_preamble.tex.
%
% Same shape as the 2026 Northwestern workshop theme
% (Teaching/PanelData/_latex/2026/202608_Northwestern/workshop_theme.tex),
% with the workshop's purple swapped for the course's cardinal. Kept
% deliberately light so the lecture sources stay easy to edit.
%
% To retune the course to Stanford's official cardinal, change the one
% Cardinal line below to HTML 8C1515.
% =====================================================================
\definecolor{Cardinal}{HTML}{A8003B}     % primary: titles, structure, alerts
\definecolor{CourseInk}{HTML}{252A34}    % body text
\definecolor{CourseBlue}{HTML}{2F6B9A}   % secondary accent
\definecolor{CourseMist}{HTML}{F8F1F4}   % block backgrounds

\setbeamercolor{normal text}{fg=CourseInk,bg=white}
\setbeamercolor{structure}{fg=Cardinal}
\setbeamercolor{alerted text}{fg=Cardinal}
\setbeamercolor{frametitle}{fg=white,bg=Cardinal}
\setbeamercolor{title}{fg=Cardinal}
\setbeamercolor{subtitle}{fg=CourseInk}
\setbeamercolor{block title}{fg=white,bg=Cardinal}
\setbeamercolor{block body}{fg=CourseInk,bg=CourseMist}
\setbeamercolor{section in toc}{fg=Cardinal}
\setbeamercolor{subsection in toc}{fg=CourseInk}
\setbeamercolor{itemize item}{fg=Cardinal}
\setbeamercolor{itemize subitem}{fg=Cardinal}
\setbeamercolor{enumerate item}{fg=Cardinal}

\setbeamerfont{title}{series=\bfseries,size=\Large}
\setbeamerfont{frametitle}{series=\bfseries,size=\large}
\setbeamerfont{block title}{series=\bfseries}
\setbeamertemplate{navigation symbols}{}
\setbeamertemplate{items}[circle]
\setbeamertemplate{blocks}[rounded][shadow=false]
\setbeamertemplate{frametitle continuation}{}
\setbeamertemplate{footline}{%
  \leavevmode\hbox{%
  \begin{beamercolorbox}[wd=.88\paperwidth,ht=2.6ex,dp=1.2ex,leftskip=1.2em]{author in head/foot}%
  \usebeamerfont{author in head/foot}\insertshorttitle
  \end{beamercolorbox}%
  \begin{beamercolorbox}[wd=.12\paperwidth,ht=2.6ex,dp=1.2ex,center]{date in head/foot}%
  \insertframenumber
  \end{beamercolorbox}}}

\newcommand{\slidesource}[1]{\vfill{\scriptsize\color{gray}Source: #1}}
\newcommand{\keyquestion}[1]{\begin{block}{Question}#1\end{block}}
      % colors and slide furniture

\graphicspath{{./figs/}}

\title[Multiple Regression Mechanics]{450A Political Methodology I}
\subtitle{Lecture 6: Multiple Regression \\ 1. Mechanics}
\author{Stanford University}
\date{}

% This lecture uses sections only; no subsection roadmap frames.
\AtBeginSubsection[]{}

\begin{document}
\begin{frame}[plain,c]
  \titlepage
  \addtocounter{framenumber}{-1}
\end{frame}



%%%%%%%%%%%%%%%%%

\section{Preparation}


% \frame{
% \frametitle{Matrix Algebra}


% Why do we study matrix algebra?

% \begin{itemize}

% \item We will use matrix algebra to derive the linear regression model \pause
% \item Matrices are an intuitive way to think about data. The data set is an array of numbers where: \pause
% \begin{itemize}
%   \item every row represents the measurements of multiple characteristics for a single unit (e.g. age and gender for person $i$)
%   \item every column represents the measurements for a single characteristic across units (e.g. age for person $i=1,...,n$)
% \end{itemize}\pause
% \item Matrices allow us to solve linear systems of equations \pause
% \item Notation is more compact and concise so many articles use it \pause
% \item Very useful for advanced methods and statistical computing
% \end{itemize}

% }

% \section{Matrix Algebra Review}

% \subsection{Matrices and Vectors}

% \frame{
% \frametitle{Definition of Matrices and Vectors}

% A \alert{matrix} is simply an arrangement of numbers in rectangular form. Generally, a
%  \( (j \times k) \) matrix \( \mathbf{A} \) can be written as follows:

%   \[ \mathbf{A} = \left[
% \begin{array}{cccc}
% 	 a_{1,1} & a_{1,2} & \cdots & a_{1,k}  \\
% 	 a_{2,1} & a_{2,2} & \cdots & a_{2,k}  \\
% 	 \vdots & \vdots & \ddots & \vdots  \\
% 	 a_{j,1} & a_{j,2} & \cdots & a_{j,k}
% \end{array}
% \right] \]

% \begin{itemize}
% \item Note that the \alert{elements} are double sub-scripted, with the row number first, and the column number second
% \item The\alert{ dimensionality} (or \alert{order}) of this matrix is \alert{$j \times k$},  where $j$ is the number of rows and $k$ is the number of columns.
% %\item The dimensionality of the matrix is also called the \alert{order} of a matrix. In general terms, the \( \mathbf{A} \) above is of order $(j, k)$.
% \item When reading the text, and doing your assignments, you should always keep track of the dimensionality.
% \end{itemize}
% }

% \frame{
%  \frametitle{A Few Examples}
% The matrix
% \[ \mathbf{W} = \left[
% \begin{array}{cc}
% 	1  & 3  \\
% 	 2 & -6
% \end{array}
% \right] \]
% is $2 \times 2$. This is also called a \alert{square matrix}, because the row dimension
% equals the column dimension $(j = k)$.\\\bigskip There are also\alert{ rectangular matrices}
% $(j \neq k)$, such as:
% \pause
% \[ \boldsymbol{G} = \left[
% \begin{array}{cc}
% 	 1 & 4  \\
% 	 1 & 3  \\
% 	 1 & -2  \\
% 	 0 & 3
% \end{array}
% \right] \]
% which is $4 \times 2$.
% \vspace{.5cm}
% \pause
% Matrices are often written with bold, CAPITAL letters.\\ Roman ($\boldsymbol{X} $) is typically data, and Greek ($\boldsymbol{\beta} $) is, as before, parameters.

%  }

% \frame{
% \frametitle{Vectors}
% Vectors are special matrices that have either \alert{one row} or \alert{one column}.\\
% \vspace{.4cm}

%  \alert{Row vectors} are $1 \times k$. For example:
% \[ \boldsymbol{\alpha}=\left[
% \begin{array}{ccccc}
% 	\alpha_{1}  & \alpha_{2} & \alpha_{3} & \cdots & \alpha_{k}
% \end{array}
% \right] \] \\A $j \times k$ matrix is a collection of $j$ $1 \times k$  row vectors.\\\bigskip

% \pause
% \alert{Column vectors} are $k \times 1$. For example,

% \[ \mathbf{y} = \left[
% \begin{array}{c}
% 	y_{1}   \\
% 	y_{2}  \\
% 	\vdots   \\
% 	y_{k}
% \end{array}
% \right] \]\\
% A $k \times j$ matrix is a collection of $j$ $k \times 1$  column vectors.\\\bigskip
% Vectors are usually bold, lowercase letters ($\mathbf{y}$). Unless otherwise specified, we assume that vectors are column vectors.
% }

% %	\frame{
% %	\frametitle{Matrix Addition and Subtraction}
% %	
% %	Now that we have defined matrices, vectors, and scalars (which are just single numbers), we can start to consider
% %	the operations we can perform on them.\\\bigskip Given a matrix of numbers, one can
% %	extend regular scalar algebra in a straight forward way.\\\bigskip
% %	Scalar addition is simply: $m + n = 2 + 5 = 7$\\\bigskip
% %	
% %	Addition is similarly defined for matrices: If matrices or vectors are of the same order, then they can be added. One
% %	performs the addition \alert{element by element}.
% %	
% %	}

% \subsection{Matrix Operations: Addition and Subtraction}

% \frame{
% \frametitle{Matrix Addition and Subtraction}

% If matrices or vectors are of the same \alert{order}, then they can be added. One
% performs the addition \alert{element by element}.\\\bigskip

% Thus, for a pair of  $(2 \times 2)$ matrices, addition proceeds as follows for the
% problem $\mathbf{A} + \mathbf{B} = \mathbf{C}$:

% \begin{small}
% \[
% \left[
%   \begin{array}{cc}
%     a_{1,1} & a_{1,2} \\
%     a_{2,1} & a_{2,2} \\
%   \end{array}
% \right] + \left[
%   \begin{array}{cc}
%     b_{1,1} & b_{1,2} \\
%     b_{2,1} & b_{2,2} \\
%   \end{array}
% \right]= \left[
%   \begin{array}{cc}
%     a_{1,1} + b_{1,1} & a_{1,2} + b_{1,2} \\
%     a_{2,1} + b_{2,1} & a_{2,2} + b_{2,2} \\
%   \end{array}
% \right] = \left[
%   \begin{array}{cc}
%     c_{1,1} & c_{1,2} \\
%     c_{2,1} & c_{2,2} \\
%   \end{array}
% \right]
% \]
% \end{small}

% Subtraction similarly follows.\\\bigskip

% Matrix addition and subtraction is only defined for matrices that share the same dimensionality. If they do, they are said to be \alert{conformable}
% for addition. If not, they are \alert{nonconformable} for addition.

% }

% \frame{
% \frametitle{Matrix Addition and Subtraction}

% Subtraction Example:
% \[
% \left[
% \begin{array}{ccc}
%     1 & 4 & -2  \\
%     5 & -3 & 3
% \end{array}
% \right]
% -
% \left[
% \begin{array}{ccc}
%     -3 & 2 & 8  \\
%     2 & 2 & -3
% \end{array}
% \right]
% = \pause
% \left[
% \begin{array}{ccc}
%     4 & 2 & -10  \\
%     3 & -5 & 6
% \end{array}
% \right]
% \]
% \pause
% Is matrix addition is\alert{ commutative}?  \[
% \mathbf{A} + \mathbf{B} = \mathbf{B}+ \mathbf{A}\]
% \bigskip \pause Yes it is.
% \pause
% Is matrix addition is \alert{associative}? \[
% \mathbf{A} + (\mathbf{B} + \mathbf{C} )=(\mathbf{A} + \mathbf{B}) + \mathbf{C} \]
% \bigskip \pause Yes it is.

% }

% \subsection{Matrix Operations: Transposition}


% \frame{
% \frametitle{Transposition}

% \alert{Transposition} of a matrix exchanges the order subscripts for each element. Thus, an order $(j \times k)$ matrix becomes an $(k \times j)$ matrix.\\\bigskip

% The transpose operation is typically denoted by placing a $'$ or $\top$ after a matrix.\\\bigskip


% Here is an example:
% \[\mathbf{Q}=\left[
%   \begin{array}{cc}
%     q_{1,1} & \textcolor{green}{q_{1,2}} \\
%     q_{2,1} & q_{2,2} \\
%     q_{3,1} & q_{3,2} \\
%   \end{array}
% \right]
% \,\, \mathbf{Q'} = \left[
%   \begin{array}{ccc}
%     q_{1,1} & q_{2,1} & q_{3,1} \\
%     \textcolor{green}{q_{1,2}} & q_{2,2} & q_{3,2} \\
%   \end{array}
% \right]\]

% Note that the subscripts in the transpose remain the same, they are just
% exchanged so $q_{1,2}$ is now in the second row, first column $(2,1)$.

% }

% \frame{
% \frametitle{Transposition}

% What happens when we transpose $\mathbf{\omega}$, a $k \times 1$ column vector?

% \[
% \boldsymbol{\omega} = \left[
% \begin{array}{r}
%     1  \\
%     3  \\
%     2  \\
%     -5
% \end{array}
% \right]\,\,\,\,\,\,\,\,\,\,\,\,\,\,\,\,
% \boldsymbol{\omega}' = \pause \left[
% \begin{array}{cccc}
%     1 & 3 & 2 & -5
% \end{array}
% \right]   \]
% it turns into a  ($1 \times k$) row vector. We can switch between row and column vectors with the transpose operation.\\
% \bigskip

% For any matrix $\mathbf{Q}$ what is $(\mathbf{Q'})'$? \pause It's $\mathbf{Q}$. In words, the transpose of the transpose is the original matrix.\\\bigskip \pause

% It can be shown that for any two matrices of the same order, the transpose of the
% sum is equal to the sum of the transposes. Symbolically: $(\mathbf{A} + \mathbf{ B})' = \mathbf{A}' + \mathbf{B}'$

% }

% \frame{
% \frametitle{Transposition}

% A matrix $\mathbf{A}$ with dimensionality $(j \times j)$ is said to be \alert{symmetric} iff  $\mathbf{A}' =  \mathbf{A}$.\\\bigskip Can you come up with an example (e.g.  $3 \times 3$)?\\\pause


% \[ \boldsymbol{\omega} = \left[
%   \begin{array}{ccc}
%     1 & .2 & -.5 \\
%     .2 & 1 & .4 \\
%     -.5 & .4 & 1 \\
%   \end{array}
% \right] \,\,\,\,\,\,\,\,\,\,\,\,\,\,\,\,
% \boldsymbol{\omega}' =  \left[
%   \begin{array}{ccc}
%     1 & .2 & -.5 \\
%     .2 & 1 & .4 \\
%     -.5 & .4 & 1 \\
%   \end{array}
% \right]\]\bigskip

% Notice: All symmetric matrices are square, but not all square matrices are symmetric.

% }

% \subsection{Matrix Operations: Multiplication}


% \frame{
% \frametitle{Matrices and Multiplication}

% There are two types of multiplication for matrices:
% \begin{itemize}
% \item \alert{Scalar multiplication}: A scalar  multiplied with a matrix
% \item \alert{Matrix multiplication}: A matrix multiplied with a matrix
% \end{itemize}\bigskip
% \pause
% \alert{Scalar multiplication} is straightforward: A scalar times a matrix equals the scalar times each element of the matrix. Thus,

% \begin{small}
% \[ \mathbf{A}=\alpha \left[
%   \begin{array}{cc}
%     a_{1,1} & a_{1,2} \\
%     a_{2,1} & a_{2,2} \\
%   \end{array}
% \right] = \left[
%   \begin{array}{cc}
%   \alpha\,  a_{1,1} & \alpha\, a_{1,2} \\
%    \alpha\, a_{2,1} & \alpha\,a_{2,2} \\
%   \end{array}
% \right] \]
% \end{small}\pause\bigskip
% So,
% \begin{small}
% \[ \mathbf{A}=\frac{1}{2} \left[
%   \begin{array}{cc}
%     4 & 8 \\
%     6 & 2 \\
%   \end{array}
% \right] = \pause\left[
%   \begin{array}{cc}
%   2 & 4 \\
%   3 & 1 \\
%   \end{array}
% \right] \]
% \end{small}

% }

%  \frame{
% \frametitle{Matrix Multiplication}
% \small
% \alert{Matrix multiplication} is more involved: Given an ($m \times n$) matrix $\mathbf{A}$ and an ($n \times r$) matrix $\mathbf{B}$, we can define the product $\mathbf{AB} = \mathbf{C}$ as the  ($m \times r$) matrix whose entries are defined by:

% \[
% \sum_{k=1}^n \, a_{i,k} \, b_{k,j} \ = \ c_{i,j}
% \]
% where $i=1,...,m$ and $j=1,...,r$.\\\bigskip
% \vspace{-.5cm}
% \[ \mathbf{A} = \left[
% \begin{array}{rrr}
%     -2 & 1 & 3  \\
%     4 & 1 & 6
% \end{array}
% \right]\ \ \ \mathbf{B} = \left[
% \begin{array}{rr}
%     3 & -2  \\
%     2 & 4  \\
%     1 & -3
% \end{array}
% \right] \]
% \begin{eqnarray*}
%     \mathbf{C} & = \pause & \left[
% \begin{array}{rr}
%     -2 \cdot 3 + 1 \cdot 2 + 3 \cdot 1 & -2 \cdot -2 + 1 \cdot 4 +
%     3 \cdot -3   \\
%     4 \cdot 3 + 1 \cdot 2 + 6 \cdot 1 & 4 \cdot -2 + 1 \cdot 4 + 6
%     \cdot -3
% \end{array}
% \right]   \\
%      & = & \left[
% \begin{array}{rr}
%     -1 & -1  \\
%     20 & -22
% \end{array}
% \right]
% \end{eqnarray*}

% %Notice:
% %\begin{itemize}
% %\item For matrices to  be \alert{conformable for multiplication}, the number of columns of $\mb{A}$, $n$, must match the number of rows of $\mb{B}$ , $n$.\\
% %\item The resulting matrix  $\mathbf{C}$ has dimensionality $(m \times r)$
% %\item This is easier to see in an example...
% %\end{itemize}
% $\mathbf{A}$ is $(2 \times 3)$ and $\mathbf{B}$ is $(3 \times 2)$, so $\mathbf{C}$ is $(2 \times 2)$.
% }


% \frame{
% \frametitle{Examples: Matrix Multiplication}
% \small
% \[ \mathbf{A} = \left[
% \begin{array}{rrr}
%     -2 & 1 & 3  \\
%     4 & 1 & 6
% \end{array}
% \right]\ \ \ \mathbf{B} = \left[
% \begin{array}{rr}
%     3 & -2  \\
%     2 & 4  \\
%     1 & -3
% \end{array}
% \right] \]
% \begin{eqnarray*}
%     \mathbf{C} & =  & \left[
% \begin{array}{rr}
%     -2 \cdot 3 + 1 \cdot 2 + 3 \cdot 1 & -2 \cdot -2 + 1 \cdot 4 +
%     3 \cdot -3   \\
%     4 \cdot 3 + 1 \cdot 2 + 6 \cdot 1 & 4 \cdot -2 + 1 \cdot 4 + 6
%     \cdot -3
% \end{array}
% \right]   \\
%      & = & \left[
% \begin{array}{rr}
%     -1 & -1  \\
%     20 & -22
% \end{array}
% \right]
% \end{eqnarray*}

% %Notice:
% \begin{itemize}
% \item For two matrices $\mb{A}$ ($m \times n$) and $\mb{B}$ ($n \times r$) to be \alert{conformable for multiplication}, $\mb{AB}$, the number of columns of $\mb{A}$, $n$, must match the number of rows of $\mb{B}$, $n$.
% \item The resulting matrix  $\mathbf{C}$ has dimensionality $(m \times r)$
% \item To describe the product $\mathbf{AB}$, we say that $\mathbf{B}$ is pre-multiplied by $\mathbf{A}$
% \end{itemize}


% }

% \frame{
% Can we also compute $\mb{BA}$?

% \[ \mathbf{A} = \left[
% \begin{array}{rrr}
%     -2 & 1 & 3  \\
%     4 & 1 & 6
% \end{array}
% \right]\ \ \ \mathbf{B} = \left[
% \begin{array}{rr}
%     3 & -2  \\
%     2 & 4  \\
%     1 & -3
% \end{array}
% \right] \]
%  \pause Yes, since $\mathbf{B}$ is $(3 \times 2)$ and $\mathbf{A}$ is $(2 \times 3)$ the matrices are conformable for this product as well. What is  $\mb{BA}$?
% \[ \mathbf{BA} = \pause \left[
% \begin{array}{rrr}
%     -14 & 1 & -3  \\
%     12 & 6 & 30  \\
%     -14 & -2 & -15
% \end{array}
% \right] \]
% However, notice that $\mb{BA}$ this is not the same as:
% \begin{eqnarray*}
%     \mathbf{AB} & = & \left[
% \begin{array}{rr}
%     -1 & -1  \\
%     20 & -22
% \end{array}
% \right]
% \end{eqnarray*}
% Hence, matrix multiplication is \alert{not commutative} $\mb{AB}\neq \mb{BA}$\\\bigskip

% }

% \frame{

% In general, even when $\mb{A}$ and $\mb{B}$ are conformable for multiplication as $\mb{AB}$, they may not be conformable for multiplication as $\mb{BA}$. \\\bigskip

% For example, suppose  $\mb{A}$ is a $(2 \times 3)$ matrix and $\mb{B}$ is an $(3 \times 1)$ matrix,
% \[ \mathbf{A} = \left[
% \begin{array}{rrr}
%     -2 & 1 & 3  \\
%     4 & 1 & 6
% \end{array}
% \right]\,\,\,\,\,\,\, \mathbf{B} = \left[
%   \begin{array}{c}
%     3 \\
%     1 \\
%     2 \\
%   \end{array}
% \right]\]
% what is $\mb{AB}$? \pause
% \[ \mathbf{AB} = \left[
% \begin{array}{c}
%     1   \\
%     25
% \end{array}
% \right] \]
% which is $(2 \times 1)$.\\\bigskip What is  $\mb{BA}$? \pause Not defined since the matrices are not conformable for multiplication in this order.

% }

%  \frame{
% \frametitle{Matrix Multiplication Properties}

% \begin{itemize}
%   \item Matrix multiplication is \alert{associative}:\\
%   $(\mb{AB} ) \mb{C}=\mb{A}(\mb{BC})$  \bigskip
%   \item Matrix multiplication is \alert{distributive}:\\
%   $\mb{A} (\mb{B}+\mb{C}) =  \mb{A}\mb{B} +  \mb{A}\mb{C}$\bigskip
%   \item It can be shown that:\\
%   $(\mb{AB})' = \mb{B}'\mb{A}'$
%   \end{itemize}
% }


% \frame{
% \frametitle{Matrix Multiplication and Vectors}

% Since vectors are just special cases of matrices we can use the same formula for vector multiplication. For example let us
% define the  $(n \times 1)$  column vector $\mb{\hat{u}}$ of regression residuals. What is\\
% \[
% \mb{\hat{u}}'\mb{\hat{u}} = \pause \left[
%   \begin{array}{cccc}
%     \widehat{u}_1 & \widehat{u}_2 & \cdots & \widehat{u}_n \\
%   \end{array}
% \right] \left[
%   \begin{array}{c}
%     \widehat{u}_1 \\
%     \widehat{u}_2 \\
%     \vdots \\
%     \widehat{u}_n \\
%   \end{array}
% \right]
% \] \pause
% \[ \mb{\hat{u}}'\mb{\hat{u}} =   \widehat{u}_1  \widehat{u}_1 +   \widehat{u}_2   \widehat{u}_2 + \cdots +   \widehat{u}_n   \widehat{u}_n=\pause\sum_{i=1}^n   \widehat{u}_i^2 \]
% the result of this multiplication is a scalar, the SSR.\\\bigskip

%  \pause Matrix multiplication of two vector $\mb{\hat{u}}'\mb{\hat{v}}$ is called the \alert{inner product} or \alert{dot product} of $\mb{\hat{u}}$ and $\mb{\hat{v}}$. In general: $\mb{\hat{u}}'\mb{\hat{v}}=\sum_{i=1}^n u_i\, v_i$  \\\bigskip

% %What are the dimension of the \alert{outer product} $\mb{\hat{u}}\mb{\hat{u}'}$  ? \pause A $(n \times n)$ matrix.

% }


% %\frame{
% %\frametitle{Matrix Multiplication and Vectors}
% %
% %The \alert{inner product} of two vectors $\mb{u}'\mb{v}$ has geometric definition in Euclidean space. In such a space, a vector is a geometrical object that possesses both a magnitude and a direction. The magnitude of the vector is denoted by ||\mb{u}} and the inner product is defined by  $\mb{u}' \mb{v}=$
% %
% %
% % a  called   $\mb{\hat{u}}$ and $\mb{\hat{v}}$. In general: $\mb{\hat{u}}'\mb{\hat{v}}=\sum_{i=1}^n u_i\, v_i$  \\\bigskip
% %
% %}

% \frame{
% \frametitle{Matrix Multiplication and Vectors}

% Now consider the matrix $\mathbf{Y}$ which we can define as:
% \[ \mathbf{Y} =
% \left[
%   \begin{array}{ccc}
%     y_{1,1} &  y_{1,2} &  y_{1,3} \\
%     y_{2,1} &  y_{2,2} &  y_{2,3} \\
%   \end{array}
% \right]=\left[
%                \begin{array}{c}
%                  \mathbf{y_1}' \\
%                  \mathbf{y_2}' \\
%                \end{array}
%              \right]  \]
% with row vectors $\mathbf{y_1}' = \left[
%                                          \begin{array}{ccc}
%                                             y_{1,1} &  y_{1,2} &  y_{1,3} \\
%                                          \end{array}
%                                        \right]\,\, \mathbf{y_2}' = \left[
%                                          \begin{array}{ccc}
%                                              y_{2,1} &  y_{2,2} &  y_{2,3} \\
%                                          \end{array}
%                                        \right]$\\\bigskip \pause
% Similarly, we can define
% \[ \mathbf{X} =
% \left[
%   \begin{array}{cc}
%     x_{1,1} & x_{1,2} \\
%     x_{2,1} & x_{2,2} \\
%     x_{3,1} & x_{3,2} \\
%   \end{array}
% \right] = \left[
%             \begin{array}{cc}
%              \mathbf{x_1}  & \mathbf{x_2} \\
%             \end{array}
%           \right]\]
% where $\mathbf{x_1}$ and $\mathbf{x_2}$ represent the columns of $\mathbf{X}$.
% }

% \frame{
% \frametitle{Matrix Multiplication and Vectors}
% Given
% \[ \mathbf{Y} =
% \left[
%   \begin{array}{ccc}
%     y_{1,1} &  y_{1,2} &  y_{1,3} \\
%     y_{2,1} &  y_{2,2} &  y_{2,3} \\
%   \end{array}
% \right]=\left[
%                \begin{array}{c}
%                  \mathbf{y_1}' \\
%                  \mathbf{y_2}' \\
%                \end{array}
%              \right]
% \]
% \[
%   \mathbf{X} =
% \left[
%   \begin{array}{cc}
%     x_{1,1} & x_{1,2} \\
%     x_{2,1} & x_{2,2} \\
%     x_{3,1} & x_{3,2} \\
%   \end{array}
% \right] = \left[
%             \begin{array}{cc}
%              \mathbf{x_1}  & \mathbf{x_2} \\
%             \end{array}
%           \right]\]

% the product of the matrices, then, can be expressed in terms of four inner products.

% \[ \mathbf{YX} = \mathbf{C} = \left[
%                                    \begin{array}{cc}
%                                      \mathbf{y_1}'\mathbf{x_1} & \mathbf{y_1}'\mathbf{x_2} \\
%                                      \mathbf{y_2}'\mathbf{x_1} & \mathbf{y_2}'\mathbf{x_2} \\
%                                    \end{array}
%                                  \right] \]
% %This is the same as the summation definition of multiplication.

% }

% \frame{
% \frametitle{Special Vector and Matrix Multiplications}

% A special vector is $\mb{1}$ which is just a $(n \times 1)$ column vector of ones. This is useful because given another $(n \times 1)$  column vector of data $\mb{x}$ we can compute the inner product:\\
% \[
% \mb{1}'\mb{x} = \pause (1\cdot x_1 + 1\cdot x_2 + ... +  1\cdot x_n) = \sum_{i=1}^n x_i \]
% And also:\[
% \frac{1}{n}\mb{1}'\mb{x}=\frac{1}{n}\sum_{i=1}^n x_i = \bar{x}
% \]
% \pause
% That is, we can compute the sum of all the elements of vector $\mb{x}$ by pre-multiplying
% it by $\mb{1}'$.

% }




% \frame{
% \frametitle{Special Vector and Matrix Multiplications}
% \small
% Two vectors $\mb{x}$ and $\mb{y}$ are said to be \alert{orthogonal} (i.e. they are perpendicular) if
% and only if their inner product is equal to zero:
% $$
% \mb{x}'\mb{y} = \mb{y}'\mb{x} = \sum_{i=1}^n x_{i} y_i = 0
% $$
% \pause
% If $\mb{x}$ and $\mb{y}$ are orthogonal we write $\mb{x} \perp
% \mb{y}$. For example, if
%  \[ \mathbf{y} = \left[
% \begin{array}{c}
% 	y_{1}   \\
% 	y_{2}
% \end{array}
% \right] = \left[
% \begin{array}{c}
% 	1   \\
% 	1
% \end{array}
% \right] \]
% \[ \mathbf{x} = \left[
% \begin{array}{c}
% 	x_{1}   \\
% 	x_{2}
% \end{array}
% \right] = \left[
% \begin{array}{c}
% 	-1  \\
% 	1
% \end{array}
% \right] \]  The inner product is:
% $$
% \mb{x}'\mb{y} = \pause (1\cdot 1) + (-1\cdot 1) = 0
% $$

% \pause

% Note: Orthogonality implies that the the vectors are linearly independent, but linear independence does not imply orthogonality.

% %is distinct from, though closely related to, independence.
% %Two vectors can be orthogonal but not independent, and vice versa.
% }




% \subsection{Identity Matrix}

% \frame{
% \frametitle{Special Vector and Matrix Multiplications}

% When performing scalar algebra, we know that $x \cdot 1 = x$, which is known as the identity relationship. There is a similar relationship in matrix algebra:\[ \mb{A}\mb{I} = \mb{A}\] where $\mb{I}$ is called the \alert{identity matrix}.\\\bigskip

% $\mb{I}$ is a \alert{diagonal}, \alert{square} matrix with ones on the main
% diagonal, and zeros on the off diagonal.\\\bigskip

% For example, the order three identity matrix is:

% \[\mb{I}=\left[
%   \begin{array}{ccc}
%     1 & 0 & 0 \\
%     0 & 1 & 0 \\
%     0 & 0 & 1 \\
%   \end{array}
% \right]\]


% }


% \frame{
% \frametitle{Special Vector and Matrix Multiplications}

% Notice that $\mb{I}$ is oftentimes subscripted to denote its dimensionality
% (e.g. $\mb{I}_3$). Here is an
% example of the use of an identity matrix:

% \[ \left[
%      \begin{array}{ccc}
%        7 & 2 & 3 \\
%        0 & 1 & 1 \\
%        2 & -2 & 99 \\
%      \end{array}
%    \right]
% \left[
%   \begin{array}{ccc}
%     1 & 0 & 0 \\
%     0 & 1 & 0 \\
%     0 & 0 & 1 \\
%   \end{array}
% \right]= \left[
%      \begin{array}{ccc}
%        7 & 2 & 3 \\
%        0 & 1 & 1 \\
%        2 & -2 & 99 \\
%      \end{array}
%    \right]\]
% \pause
% One of the nice properties of the identity matrix is that it is \alert{commutative},
% \alert{associative}, and \alert{distributive} with respect to multiplication. That is,
% \[ \mb{A}\mb{I}\mb{B}=\mb{I}\mb{A}\mb{B}=\mb{A}\mb{B}\mb{I}=\mb{A}\mb{B}
% \]

% }

%\section{Regression in Matrix Form}

\subsection{Matrix Representation of Regression}

\frame{
\frametitle{Simple regression in Vector Form}

If we define the following column vectors:
\begin{itemize}
\item $\mathbf{1}$ denotes a vector with $n$ ones
\item $\mathbf{y}$ denotes an $n$ element data vector for $y$
\item $\mathbf{x}$ denotes an $n$ element data vector for $x$
%\item $\mathbf{\epsilon}$ denotes an $n$ element vector of errors
%\item $\mathbf{y^{*}}$ denotes an $n$ element data vector of the mean-deviations for $y$
%\item $\mathbf{x^{*}}$ denotes an $n$ element data vector of the mean-deviations for $x$
\item $\hat{\mathbf{u}}$ denotes an $n$ element vector of residuals
\item $\hat{\mathbf{y}}$ denotes an $n$ element vector of fitted values
\end{itemize}
\vspace{.5cm}
%Then we can write the simple regression model
%\[ \mb{y} = \beta_0\mb{1} + \beta_1\mb{x} + \mb{\epsilon}  \]
Then we can write the fitted regression line for the data in the following manner.
\[ \mb{y} = \hat{\beta}_0\mb{1} + \hat{\beta}_1\mb{x} + \hat{\mb{u}}  \]
and ...
\[ \hat{\mb{y}} = \hat{\beta}_0\mb{1} + \hat{\beta}_1\mb{x}   \]
}

\frame{
\frametitle{Example: Simple Regression with Three Observations}

\[ \mathbf{y} = \left[
\begin{array}{c}
	y_{1}   \\
	y_{2}  \\
	y_3
\end{array}
\right] = \left[
\begin{array}{c}
	1   \\
	2  \\
	3
\end{array}
\right]
\, ; \, \mathbf{1} = \left[
\begin{array}{c}
	1   \\
	1 \\
	1
\end{array}
\right] \, ; \, \mathbf{x} = \left[
\begin{array}{c}
	x_{1}   \\
	x_{2} \\
	x_3
\end{array}
\right] = \left[
\begin{array}{c}
	1   \\
	3  \\
	2
\end{array}
\right] \]
\bigskip
\[
 \mb{y} = \hat{\beta}_0\mb{1} + \hat{\beta}_1\mb{x} +\hat{\mb{u}}
\]\bigskip
\[
  \left[
\begin{array}{c}
	1   \\
	2  \\
	3
\end{array}
\right]  =  \hat{\beta}_0 \left[
\begin{array}{c}
	1   \\
	1 \\
	1
\end{array}
\right] + \hat{\beta}_1 \left[
\begin{array}{c}
	1   \\
	3  \\
	2
\end{array}
\right] +  \left[
\begin{array}{c}
	\hat{u}_1  \\
	\hat{u}_2  \\
	\hat{u}_3
\end{array}
\right]
\]

}

\frame{
\frametitle{Multiple Regression in Matrix Form}
\small
\begin{small}
Let $\mb{X}$ be our $n \times (k + 1)$ data matrix for $k$ predictors (plus the constant),
\[ \mathbf{X} = \left[
\begin{array}{ccccc}
    1 & x_{11} &  x_{12} & \dots & x_{1k}  \\
    1 & x_{21} & x_{22}  & \dots & x_{2k} \\
   \vdots & \vdots & \vdots & \vdots & \vdots \\
   1 & x_{n1} & x_{n2}  & \dots & x_{nk}
\end{array}
\right] \]
\pause
Let $\hat{\boldsymbol{\beta}}$ be the vector $(k + 1) \times 1$ of least squares estimates (intercept and $k$ slopes): $\hat{\boldsymbol{\beta}} = \left[
\begin{array}{c}
    \hat{\beta}_0 \\
    \hat{\beta}_1 \\
    \vdots \\
    \hat{\beta}_k
    \end{array}
\right]$\\
\pause
Then $\mb{X}\hat{\boldsymbol{\beta}}$ is the $n \times 1$ vector of fitted values:
\[
\hat{\mb{y}} = \left[
\begin{array}{c}
    \hat{y}_1 \\
    \hat{y}_2 \\
    \vdots \\
    \hat{y}_n
    \end{array}
\right]  = \pause \mb{X}\hat{\boldsymbol{\beta}} = \pause
\left[
\begin{array}{c}
    1\hat{\beta}_0  + x_{11}\hat{\beta}_1 +  x_{12}\hat{\beta}_2 + \dots + x_{1k} \hat{\beta}_k \\
    1\hat{\beta}_0 + x_{21}\hat{\beta}_1 + x_{22}\hat{\beta}_2  + \dots + x_{2k}\hat{\beta}_k \\
   \vdots  \\
   1\hat{\beta}_0 + x_{n1}\hat{\beta}_1 + x_{n2}\hat{\beta}_2  + \dots + x_{nk}\hat{\beta}_k
\end{array}
\right]
\]
\end{small}
}

\begin{frame}[fragile]
\frametitle{Explanatory Data Matrix}
The $n \times (k + 1)$ matrix $\mb{X}$ is often called the \alert{design matrix}. Example:

% Output of code/fig_chile_matrices.R (saved in code/fig_chile_matrices.txt).
\begin{Verbatim}[fontsize=\footnotesize, frame=single, label=R Code]
> fit <- lm(vote1 ~ fem + educ + age, data = d)
> X <- model.matrix(fit)
> dim(X)
[1] 1703    4
> head(X, 8)
  (Intercept) fem educ age
1           1   0    1  65
2           1   0    3  29
3           1   1    1  38
4           1   1    1  49
5           1   1    2  23
6           1   1    1  28
7           1   0    3  26
8           1   1    2  24
\end{Verbatim}
In R, the design matrix can be extracted from a fitted {\tt lm} object
by the {\tt model.matrix} function.
\slidesource{national survey by FLACSO/Chile, April--May 1988, before the plebiscite (\texttt{carData::Chile}); we keep the 1,703 of 2,700 respondents who intended to vote Yes or No and answered every question. \texttt{vote1} $=1$ if the respondent intends to vote Yes (for Pinochet), $0$ if No; \texttt{fem} $=1$ if female; \texttt{educ}: 1 = primary, 2 = secondary, 3 = post-secondary.\par}
\end{frame}

\subsection{Solving Linear Systems using Matrices}

\frame{
\frametitle{Solving Systems of Equations Using Matrices}

Matrices are very useful to solve linear systems of equations, such as the first order conditions for our least squares estimates.\\\bigskip

Here is an example with three equations and three unknowns:

\begin{eqnarray*}
    x + 2y + z & = & 3  \\
    3x - y - 3z & = & -1  \\
    2x + 3y + z & = & 4
\end{eqnarray*}

How would one go about solving this? \pause There are various techniques, including
substitution, and multiplying equations by constants and adding them to get
single variables to cancel.

}

\frame{
\frametitle{Solving Systems of Equations Using Matrices}
\small
An easier way is to use matrix algebra. Note that the system of equations
\begin{eqnarray*}
    x + 2y + z & = & 3  \\
    3x - y - 3z & = & -1  \\
    2x + 3y + z & = & 4
\end{eqnarray*}
can be written as follows:
\[
\left[
\begin{array}{rrr}
    1 & 2 & 1  \\
    3 & -1 & -3  \\
    2 & 3 & 1
\end{array}
\right]
\left[
\begin{array}{r}
    x  \\
    y  \\
    z
\end{array}
\right]
=
\left[
\begin{array}{r}
    3  \\
    -1 \\
    4
\end{array}
\right] \iff
\mathbf{Au = b}
\]

How do we solve this for $\mb{u}=\left[
\begin{array}{r}
    x  \\
    y  \\
    z
\end{array}
\right]$?  Let's look at the scalar case first.

}




\frame{
\frametitle{Solving Equations with Inverses (scalar case)}
\small
Let's go back to the scalar world of 8th grade algebra. How would you solve the following for $u$?
$$
a u = b
$$
\pause We multiply both sides of by the reciprocal $1/a$ (the \alert{inverse of $a$}) and get:
\begin{align*}
\frac{1}{a} a\, u = \frac{1}{a} b \\
u = \frac{b}{a}
\end{align*}

(Note that this technique only works if $a \ne 0$.  If $a=0$, then there are either an infinite number of solutions for $u$ (when $b=0$), or no solutions for $u$ (when $b \ne 0$).)\\\bigskip
\pause
So to solve our multiple equation problem in the matrix case we need a matrix equivalent of the inverse. This equivalent is the \alert{inverse matrix}. The inverse of $\mb{A}$ is written as $\mb{A^{-1}}$.

}

\frame{
\frametitle{Inverse of a Matrix}
\small
The \alert{inverse} $\mb{A}^{-1}$ of $\mb{A}$ has the property that $\mb{A}^{-1}\mb{A} =
\mb{A}\mb{A}^{-1} = \mb{I}$ where $\mb{I}$ is the identity matrix.

\begin{itemize}
  \item The inverse $\mb{A}^{-1}$ exists only if $\mb{A}$ is invertible or \alert{nonsingular} (more on this soon)
%  \begin{itemize}
%    \item All invertible matrices are square, but not all square matrices are invertible.
%    \item $\mb{A}$ is invertible iff it is \alert{nonsingular}. We can check singularity of $\mb{A}$ in a number of ways (e.g. computing the column rank).
%  \end{itemize}
  \item The inverse is unique if it exists and then the linear system has a unique solution.
  \item There are various methods for finding/computing the inverse of a matrix
\end{itemize}

The \alert{inverse matrix} allows us to solve linear systems of equations.

\begin{align*}
\uncover<2->{\mathbf{Au = b}\\}
\uncover<3->{\mathbf{A^{-1}Au = A^{-1}b}\\}
\uncover<4->{\mathbf{Iu = A^{-1}b}\\}
\uncover<5->{\mathbf{u = A^{-1}b}}
\end{align*}

}


\frame{
\small
Given $\mb{A}$ we find that $\mb{A^{-1}}$ is:
\[\mb{A}=
\left[
\begin{array}{rrr}
    1 & 2 & 1  \\
    3 & -1 & -3  \\
    2 & 3 & 1
\end{array}
\right]\,;\, \mb{A^{-1}}=
\left[
\begin{array}{rrr}
    8 & 1 & -5  \\
    -9 & -1 & 6  \\
    11 & 1 & -7
\end{array}
\right]\]
We can now solve our system of equations:\pause
\[
\mb{u}=\mb{A^{-1}b}=
\left[
\begin{array}{rrr}
    8 & 1 & -5  \\
    -9 & -1 & 6  \\
    11 & 1 & -7
\end{array}
\right]\left[
\begin{array}{r}
    3  \\
    -1 \\
    4
\end{array}
\right]=\left[
\begin{array}{r}
    3  \\
    -2 \\
    4
\end{array}
\right]\]\pause
So the solution vector is $\alert{x=3}$, $\alert{y=-2}$, and $\alert{z=4}$. Verifying:
\begin{eqnarray*}
    x + 2y + z  &= \alert{3} + 2\cdot \alert{-2} + \alert{4} & =  3 \\
    3x - y - 3z  &= 3\cdot \alert{3} - \alert{-2} - 3\cdot \alert{4}& =  -1\\
    2x + 3y + z &=  2\cdot \alert{3} + 3 \cdot \alert{-2} + \alert{4} & =  4
\end{eqnarray*}
Computationally, this method is very convenient. We ``just'' compute the inverse, and perform a single matrix multiplication.
}

\frame{
\frametitle{Singularity of a Matrix}

\small

If the inverse of $\mb{A}$ exists, then the linear system has a unique solution. If it exists, we say that $\mb{A}$ is \alert{nonsingular} or \alert{invertible} (these statements are equivalent).\\\bigskip

$\mb{A}$ must be square to be invertible, but not all square matrices are invertible. More precisely, a square matrix $\mb{A}$ is invertible iff its column vectors (or equivalently its row vectors) are \alert{linearly independent}.\\\bigskip

The column \alert{rank} of a matrix $\mb{A}$ is the largest number of linearly independent columns of $\mb{A}$. If the rank of $\mb{A}$ equals the number of columns of $\mb{A}$, then we say that $\mb{A}$ has \alert{full column rank}. This implies that all its column vectors are linearly independent.\\\bigskip

%The column rank of a matrix A is the maximal number of linearly independent columns of A. Likewise, the row rank is the maximal number of linearly independent rows of A.
% Since the column rank and the row rank are always equal, they are simply called the rank of A.

If a column of $\mb{A}$ is a linear combination of the other columns, there are either \alert{no solutions} to the system of equations or infinitely \alert{many solutions} to the system of equations. With no solution, the system is called \alert{inconsistent}; with infinitely many, the equations do not pin down a single solution.

}

%\frame{
%\frametitle{Geometric Example in 2D}
%\begin{figure}[ht] \centering
%        \scalebox{.4}{\includegraphics{pres1.pdf}}
%\end{figure}
%
%\[
%A =\left[
%\begin{array}{rr}
%    1 & -1\\
%    3 &  1
%\end{array}
%\right] \, \quad \,
%A =\left[
%\begin{array}{rr}
%    4 & 3\\
%    1 & -2 \\
%    3 & 5
%\end{array}
%\right] \, \quad \,
%A =\left[
%\begin{array}{rr}
%    3 & 2\\
%    3 & 2
%\end{array}
%\right]
%\]
%}

\frame{
\frametitle{Geometric Example in 2D}
\centering
% Drawn in TikZ (2026); replaces three cropped 2015 screenshots (now in _old/figs/eq1-3.pdf).
\begin{tikzpicture}[x=0.64cm, y=0.5cm, font=\footnotesize, >=Stealth]
  % three identical panels: axes and titles
  \foreach \k/\ttl in {0/{One solution}, 1/{No solution}, 2/{Infinitely many solutions}} {
    \begin{scope}[xshift=\k*5.15cm]
      \draw[->, gray!80] (-0.6,0) -- (5.4,0) node[right, text=CourseInk] {$x$};
      \draw[->, gray!80] (0,-2.8) -- (0,5.0) node[above, text=CourseInk] {$y$};
      \foreach \t in {1,2,3,4}
        % inner sep=1pt: the labels sit close to the axis, clear of the red line in panel 1
        \draw[gray!80] (\t,0.12) -- (\t,-0.12) node[below, inner sep=1pt, font=\scriptsize, text=CourseInk] {$\t$};
      \foreach \t in {-2,2,4}
        \draw[gray!80] (0.12,\t) -- (-0.12,\t) node[left, font=\scriptsize, text=CourseInk] {$\t$};
      \node[font=\small\bfseries, text=Cardinal] at (2.4,6.2) {\ttl};
    \end{scope}
  }
  % panel 1: x + y = 4 and 2x - 3y = 6 cross once, at (3.6, 0.4)
  \begin{scope}
    \draw[CourseBlue, very thick] (-0.5,4.5) -- (5.2,-1.2);
    % the red line starts at its y-intercept (0,-2), so it stays clear of the tick label -2
    \draw[Cardinal, very thick] (0,-2) -- (5.2,1.4667);
    \node[CourseBlue, anchor=south west] at (0.55,3.45) {$x + y = 4$};
    \node[Cardinal, anchor=north west] at (0.9,-1.75) {$2x - 3y = 6$};
    \fill[CourseInk] (3.6,0.4) circle (2.2pt);
    \draw[->, CourseInk, thin] (3.6,2.55) -- (3.6,0.75);
    \node[CourseInk, anchor=south] at (3.6,2.5) {$(3.6,\ 0.4)$};
  \end{scope}
  % panel 2: x + y = 4 and x + y = 2 are parallel
  \begin{scope}[xshift=5.15cm]
    \draw[CourseBlue, very thick] (-0.5,4.5) -- (5.2,-1.2);
    \draw[Cardinal, very thick] (-0.5,2.5) -- (4.7,-2.7);
    \node[CourseBlue, anchor=south west] at (2.2,1.9) {$x + y = 4$};
    \node[Cardinal, anchor=north west] at (0.2,-1.0) {$x + y = 2$};
  \end{scope}
  % panel 3: x + y = 4 and 2x + 2y = 8 are the same line
  \begin{scope}[xshift=10.3cm]
    % same width as the other lines; the red dashes alternate with the blue line underneath
    \draw[CourseBlue, very thick] (-0.5,4.5) -- (5.2,-1.2);
    \draw[Cardinal, very thick, dash pattern=on 4pt off 3pt] (-0.5,4.5) -- (5.2,-1.2);
    \node[CourseBlue, anchor=south west] at (2.2,1.9) {$x + y = 4$};
    \node[Cardinal, anchor=south west] at (2.2,2.95) {$2x + 2y = 8$};
  \end{scope}
  % the same systems written as A u = b (panel centers are 8.047 x-units apart)
  \node at (2.4,-4.45)
    {$\mb{A}=\left[\begin{array}{rr} 1 & 1\\ 2 & -3\end{array}\right],\ \mb{b}=\left[\begin{array}{r} 4\\ 6\end{array}\right]$};
  \node at (10.447,-4.45)
    {$\mb{A}=\left[\begin{array}{rr} 1 & 1\\ 1 & 1\end{array}\right],\ \mb{b}=\left[\begin{array}{r} 4\\ 2\end{array}\right]$};
  \node at (18.494,-4.45)
    {$\mb{A}=\left[\begin{array}{rr} 1 & 1\\ 2 & 2\end{array}\right],\ \mb{b}=\left[\begin{array}{r} 4\\ 8\end{array}\right]$};
  \node at (2.4,-6.2) {$\mb{A}$ invertible};
  \node at (10.447,-6.2) {$\mb{A}$ singular};
  \node at (18.494,-6.2) {$\mb{A}$ singular};
\end{tikzpicture}
}





%\frame{
%\frametitle{Geometric Example in 2D}
%\begin{figure}[ht] \centering
%        \scalebox{.6}{\includegraphics{g1.pdf}}
%\end{figure}
%
%}

%\frame{
%\frametitle{Geometric Example in 3D}
%\begin{figure}[ht] \centering
%        \scalebox{.45}{\includegraphics{g2.pdf}}
%\end{figure}
%
%}

\frame{
\frametitle{Why do we care about invertibility?}
\begin{small}
We have seen that OLS regression is defined by a system of linear equations

\[
\hat{\mb{y}} = \left[
\begin{array}{c}
    \hat{y}_1 \\
    \hat{y}_2 \\
    \vdots \\
    \hat{y}_n
    \end{array}
\right]  = \mb{X}\hat{\boldsymbol{\beta}} =
\left[
\begin{array}{c}
    1\hat{\beta}_0  + x_{11}\hat{\beta}_1 +  x_{12}\hat{\beta}_2 + \dots + x_{1k} \hat{\beta}_k \\
    1\hat{\beta}_0 + x_{21}\hat{\beta}_1 + x_{22}\hat{\beta}_2  + \dots + x_{2k}\hat{\beta}_k \\
   \vdots  \\
   1\hat{\beta}_0 + x_{n1}\hat{\beta}_1 + x_{n2}\hat{\beta}_2  + \dots + x_{nk}\hat{\beta}_k
\end{array}
\right]
\]
with our data matrix\[ \mathbf{X} = \left[
\begin{array}{ccccc}
    1 & x_{11} &  x_{12} & \dots & x_{1k}  \\
    1 & x_{21} & x_{22}  & \dots & x_{2k} \\
   \vdots & \vdots & \vdots & \vdots & \vdots \\
   1 & x_{n1} & x_{n2}  & \dots & x_{nk}
\end{array}
\right] \]
\pause
We have also learned that $\hat{\boldsymbol{\beta}}$ is obtained by solving \alert{normal equations},
a linear system of equations. \\
\pause
It turns out that to solve for $\hat{\boldsymbol{\beta}}$, we need to invert $\mb{X'X}$, a $(k + 1) \times (k + 1) $ matrix.
\end{small}

}

\frame{
\frametitle{Some Non-invertible Explanatory Data Matrices}

\small
$\mb{X'X}$ is invertible iff $\mb{X}$ is full column rank (see Blackwell, ch.~7.4), so the collection of predictors need to be linearly independent (no perfect collinearity).\\\bigskip Some example of $\mb{X}$ that are not full column rank:
\begin{small}
\[ \mathbf{X} = \left[
\begin{array}{ccc}
    1 & 2 &  -2   \\
    1 & 3 &  -3   \\
    1 & 4 &  -4   \\
    1 & 5 &  -5
\end{array}
\right] \]
%\[ \mathbf{X} = \left[
%\begin{array}{cccc}
%    1 & 27 &  12 & 0     \\
%    1 & 46 & 16  & 1   \\
%   1 & 33 & 12  & 1
%\end{array}
%\right] \]\pause
\[ \mathbf{X} = \left[
\begin{array}{ccc}
    1 & 54 &  54{,}000      \\
    1 & 37 &  37{,}000   \\
   1 & 89 &  89{,}000   \\
   1& 72 & 72{,}000
\end{array}
\right] \]
\[ \mathbf{X} = \left[
\begin{array}{ccc}
    1 & 0 &  1      \\
    1 & 1 &  0   \\
   1 & 1 &  0   \\
   1& 0 & 1
\end{array}
\right] \]
\end{small}
}

\section{The OLS Estimator in Matrix Algebra}

\subsection{Vector Derivatives}
\frame{
\frametitle{The Linear Model in Matrix Form}

Let's consider the OLS estimator. The general multiple regression model can be written as,

$$
\mb{y} = \X \b  + \mb{u}
$$

where:
\begin{itemize}
  \item $\y$ is the $n \times 1$ vector of responses
  \item $\X$ is the $n \times (k+1)$ design matrix of predictor variables plus constant
  \item $\b$ is the $(k+1) \times 1$ vector of regression coefficients (intercept and slopes)
  \item $\mb{u}$ is the $n \times 1$ vector of errors.
\end{itemize}
}

\frame{
\frametitle{Sum of Squared Residuals}
\scriptsize
Recall that we derive the OLS estimator by minimizing the sum of squared residuals (SSR) with respect to $\b$.

\medskip
\pause

Let's define the $n \times 1$ vector of residuals for a hypothetical value of $\b$, $\tilde\b$,
as $\tilde{\mb{u}}=\y - \X \tilde{\b}$.

\medskip
\pause

Then we can write SSR as,
\begin{align*}
\uncover<2->{S(\tilde{\b}, \X, \y)
   &= \tilde{\mb{u}}'\tilde{\mb{u}} \l(= \textstyle \sum_{i=1}^n \tilde{u}_i^2 \r)\\}
\uncover<3->{   &=  (\y - \X \tilde{\b})'(\y - \X \tilde{\b}) \\}
\uncover<4->{   &=  (\y' - (\X \tilde{\b})')(\y - \X \tilde{\b}) \\}
\uncover<5->{   &=  (\y' - \tilde{\b}'\X')(\y - \X \tilde{\b}) \\}
\uncover<6->{   &= \y'\y - \tilde{\b}'\X'\y - \y'\X\tilde{\b} + \tilde{\b}'\XX  \tilde{\b} \\}
\uncover<7->{   &= \y'\y - (\X \tilde{\b})' \y - \y'\X\tilde{\b} + \tilde{\b}'\XX  \tilde{\b} \\}
\uncover<8->{   &= \y'\y - \y' (\X \tilde{\b}) - \y'\X\tilde{\b} + \tilde{\b}'\XX  \tilde{\b} \\}
\uncover<9->{   &= \y'\y - 2 \y'\X\tilde{\b} + \tilde{\b}'\XX \tilde{\b}\, %\,\,(\mbox{using}\,\, (\mb{AB})' = \mb{B}'\mb{A}')
}
\end{align*}

%Recall:  $(\mb{A-B})' = \mb{A}'-\mb{B}'$ and $(\mb{AB})' = \mb{B}'\mb{A}'$)\\

\uncover<10->{ Once we have defined the OLS estimator of $\b$, how do we find the vectors of estimates $\tilde{\b}$ that minimizes the SSR?\\\bigskip}

\uncover<11->{We compute the vector of partial derivatives and solve the resulting system of equations with an inverse matrix.}

}


\frame{
\frametitle{Gradient}
\small

%A few matrix derivative rules are needed to derive the OLS estimator.\\\medskip

Let $v = v(\mb{u})$ be a \alert{scalar-valued function} $\mathbb{R}^n \rightarrow \mathbb{R}$ where
$\mb{u}$ is a $(n \times 1)$ column vector.
For example: $v(\mb{u}) = \mb{c}'\mb{u}$ where $\mb{c} = \left[
                                                              \begin{array}{c}
                                                                0 \\
                                                                1 \\
                                                                3 \\
                                                              \end{array}
                                                            \right]$ and $\mb{u} = \left[
                                                                                     \begin{array}{c}
                                                                                       u_1 \\
                                                                                       u_2 \\
                                                                                       u_3 \\
                                                                                     \end{array}
                                                                                   \right]$ \\ \smallskip
\begin{definition}[gradient]
We can define the column \alert{vector of partial derivatives}

$$
\frac{\partial v(\mb{u})}{\partial \mb{u}} =
\left[
\begin{array}{c}
    \partial v/\partial u_1 \\
    \partial v/\partial u_2 \\
    \vdots \\
    \partial v/\partial u_{n}
\end{array}
\right]
$$

This vector of partial derivatives is called the \alert{gradient}.
\end{definition}
}

\frame{
\frametitle{Vector Derivative Rule I (linear functions)}
\small
\vspace{-.25\baselineskip}
\begin{theorem}[differentiation of linear functions]
Given a linear function $v(\mb{u}) = \mb{c}'\mb{u}$ of an ($n \times 1$) vector $\mb{u}$, the derivative of $v(\mb{u})$ w.r.t. $\mb{u}$ is given by $$\frac{\partial v} { \partial \mb{u}} = \mb{c}$$
This also works when $\mb{c}$ is a matrix and therefore $v$ is a vector-valued function.
\end{theorem}
\pause
\medskip

For example, let $v(\mb{u}) = \mb{c}'\mb{u}$ where $\mb{c} = \left[
                                                              \begin{array}{c}
                                                                0 \\
                                                                1 \\
                                                                3 \\
                                                              \end{array}
                                                            \right]$ and $\mb{u} = \left[
                                                                                     \begin{array}{c}
                                                                                       u_1 \\
                                                                                       u_2 \\
                                                                                       u_3 \\
                                                                                     \end{array}
                                                                                   \right]$
, then
$$
v = \mb{c}'\mb{u} = 0\cdot u_1 + 1 \cdot u_2 + 3 \cdot u_3
$$
and
\vspace{-.5\baselineskip}
$$
\frac{\partial v}{ \partial \mb{u}} =
\left[
\begin{array}{c}
    \partial v/\partial u_1 \\
    \partial v/\partial u_2 \\
    \partial v/\partial u_3
\end{array}
\right]
= \pause
\left[
\begin{array}{r}
    0 \\
    1 \\
    3
\end{array}
\right]
= \mb{c}
$$
Hence,
\vspace{-.5\baselineskip}
$$
\frac{\partial v} { \partial \mb{u}} = \mb{c}
$$
}

%\frame{
%\frametitle{Vector Derivatives: Example I}
%\small
%If $v(\mb{u}) = \mb{c}'\mb{u}$ then $\frac{\partial v} { \partial \mb{u}} = \mb{c}$\\\bigskip
%We care about this, since for the OLS estimator we need the derivative of $v(\mb{u}) = \mb{c'}\mb{u}$ where\\\medskip
%\begin{itemize}
%  \item  $\mb{c}' = - 2 \mb{y'X}$ is of dimensionality \pause $(1 \times k+1)$. Recall $\mb{y}$ is $(n \times 1)$ and $\mb{X}$ is $(n \times k+1)$\pause \medskip
%  \item $\mb{u}=\tilde{\boldsymbol{\beta}}$ is our column vector with $k+1$ regression coefficients
%  \end{itemize}\bigskip

%So
%$$
%v = (- 2 \mb{y'X})\tilde{\boldsymbol{\beta}}
%$$
%and
%$$
%\frac{\partial v }{\partial \tilde{\boldsymbol{\beta}}} = \pause (- 2 \mb{y'X})' = \pause - 2 \mb{X'y}
%$$

%}

\frame{
\frametitle{Vector Derivative Rule II (quadratic form)}
\small
\begin{theorem}[quadratic form]
 Given a $(n \times n)$ symmetric matrix $\mb{A}$ and a scalar-valued function $v(\mb{u}) = \mb{u}'\mb{A}\mb{u}$ of $(n \times 1)$ vector $\mb{u}$, we have
$$
\frac{\partial v }{ \partial \mb{u}} = \mb{A}'\mb{u} + \mb{A}\mb{u} =
2\mb{A}\mb{u} \,\, \,%\mbox{Note: this vector is} (1 \times n)
$$
\end{theorem}
\bigskip\pause
For example, let
$
\mb{A} =
\left[
\begin{array}{cc}
  3 & 1\\
  1 & 5
\end{array}
\right] \ \mbox{and} \ \mb{u} = \left[
                                        \begin{array}{c}
                                          u_1 \\
                                          u_2 \\
                                        \end{array}
                                      \right]$. Then $v(\mb{u}) = \mb{u}'\mb{A}\mb{u}$ is equal to \pause
\begin{align*}
%v &= \mb{u}'\mb{A}\mb{u}\\
v  &=  [3 \cdot u_1 + u_2, \, \, u_1 + 5 \cdot u_2]
\left[
\begin{array}{c}
   u_1\\
   u_2
\end{array}
\right]\\
  &= 3 u_1^2 + 2 u_1 u_2 + 5 u_2 ^2
\end{align*}
and
$$
\frac{\partial v}{ \partial \mb{u}} = \pause
\left[
\begin{array}{c}
    \partial v/\partial u_1 \\
    \partial v/\partial u_2 \\
\end{array}
\right]
= \pause
\left[
\begin{array}{c}
   6 u_1 + 2 u_2 \\
   2 u_1 + 10 u_2
\end{array}
\right]
=
2 \cdot \left[
\begin{array}{c}
   3 u_1 + 1 u_2 \\
   1 u_1 + 5 u_2
\end{array}
\right]
=
2 \mb{A} \mb{u}
$$

}

%\frame{
%\frametitle{Vector Derivatives: Example II}
%\small
%Recall that when $v = \mb{u}'\mb{A}\mb{u}$ and $\mb{A}$ is symmetric, then
%$$
%\partial v / \partial \mb{u} = \mb{A}'\mb{u} + \mb{A}\mb{u} =
%2\mb{A}\mb{u}
%$$
%We care about this, since we need such a partial derivative for the OLS estimator.\\\bigskip

%Let $\mb{A} = \mb{X'X}$ be the $(k+1) \times (k+1)$ matrix that can be calculated from the matrix of explanatory variables, and let $\mb{u} = \tilde{\boldsymbol{\beta}}$, the $(k + 1 \times 1)$ vector of possible values for the intercept and slope parameters.\\\bigskip

%Then $v(\mb{u} ) = \tilde{\boldsymbol{\beta}}'\mb{X'X} \tilde{\boldsymbol{\beta}}$, and
%$$
%\frac{\partial v} { \partial \tilde{\boldsymbol{\beta}}} = \mb{A}'\mb{u} + \mb{A}\mb{u} =
%2\mb{A}\mb{u} = 2 \mb{X'X} \tilde{\boldsymbol{\beta}}
%$$
%%\scriptsize{
%%Technical aside: $\mb{X'X}$ is positive semi-definite (psd) and symmetric. A symmetric $(n \times n)$ %matrix $\mb{A}$ is psd if $\mb{u}'\mb{A}\mb{u}\geq 0$ for all $(n \times 1)$ vectors $\mb{u}$}
%%Keep this in mind.

%}

\frame{
\frametitle{Hessian}

Suppose $v$ is a scalar-valued function $v = f(\mb{u})$ of a
 $(k+1) \times 1$ column vector $\mb{u}=\left[
                  \begin{array}{cccc}
                    u_1 & u_2 & \cdots &  u_{k+1} \\
                  \end{array}
                \right]'$

\begin{definition}[Hessian]
The $(k+1) \times (k+1)$ matrix of second-order partial derivatives of $v = f(\mb{u})$ is called the \alert{Hessian matrix} and denoted
$$
\frac{\partial^2 v}{\partial \mb{u} \partial \mb{u}'}=
\left[
  \begin{array}{ccc}
   \frac{\partial^2 v}{\partial u_1\partial u_1}  & \cdots & \frac{\partial^2 v}{\partial u_1\partial u_{k+1}} \\
    \vdots & \ddots & \vdots \\
   \frac{\partial^2 v}{\partial u_{k+1}\partial u_1} & \cdots & \frac{\partial^2 v}{\partial u_{k+1}\partial u_{k+1}} \\
  \end{array}
\right]
$$
\end{definition}
Note: The Hessian is symmetric.
%We need the Hessian matrix because it will lead to the \alert{variance-covariance matrix}.

\medskip
\pause
We will use the above rules and definitions to derive the OLS estimator next.

%$$
%\frac{\partial S(\tilde{\b}, \X, \y)}{\partial \tilde{\b}} =
%-2\X'\y + 2 \XX \tilde{\b}
%$$
%
%\begin{align*}
%\frac{\partial^2 S(\tilde{\b}, \X, \y)}{\partial \tilde{\b} \partial \tilde{\b}'} &=
% 2 \XX \tilde{\b}
%\end{align*}

}


\subsection{Derivation of OLS Estimator}


\frame{
\frametitle{Derivatives with respect to $\tilde{\b}$}

%More formally,

%$$
%\hat{\b}(\X, \y) = \underset{\tilde{\b}}{\arg \min} \
%S(\tilde{\b},\X, \y)
%$$

%This is a fairly straightforward optimization problem. So, to find the
%LS estimate for given values of $X$ and $\y$ we can use calculus.

%First, find the gradient of $S(\tilde{\b}, \X, \y)$ with respect to
%$\tilde{\b}$:
\small
\begin{align*}
S(\tilde{\b}, \X, \y)
   &=  (\y - \X \tilde{\b})'(\y - \X \tilde{\b}) \\
   &= \y'\y - 2 \y'\X\tilde{\b} + \tilde{\b}'\XX \tilde{\b}
\end{align*}
$$
\frac{\partial S(\tilde{\b}, \X, \y)}{\partial \tilde{\b}} = \pause -2\X'\y + 2 \XX \tilde{\b}
$$
\begin{itemize}
\item The first term does not contain $\tilde{\b}$
\item The second term is an example of rule I from the derivative section
\item The third term is an example of rule II from the derivative section
\end{itemize}
\pause
And while we are at it the Hessian is:
$$
\frac{\partial^2 S(\tilde{\b}, \X, \y)}{\partial \tilde{\b} \partial \tilde{\b}'}=2\XX
$$


}

\frame{
\frametitle{Solving for $\hat{\b}$}

$$
\frac{\partial S(\tilde{\b}, \X, \y)}{\partial \tilde{\b}} =
-2\X'\y + 2 \XX \tilde{\b}
$$


Setting the vector of partial derivatives equal to zero and substituting $\hat{\b}$ for $\tilde{\b}$ , we can solve for the OLS estimator.
\vspace{-.5\baselineskip}
\begin{align*}
\uncover<2->{\mb{0} &= -2\X'\y + 2 \XX \hat{\b} \\}
\uncover<3->{-2 \XX \hat{\b} &= -2\X'\y \\}
\uncover<4->{\XX \hat{\b} &= \X'\y \\}
\uncover<5->{\XXinv \XX \hat{\b} &= \XXinv \X'\y \\}
\uncover<6->{\mb{I} \hat{\b} &= \XXinv \X'\y \\}
\uncover<7->{\hat{\b} &= \XXinv \X'\y }
\end{align*}
\uncover<8->{This needs $\XX$ to be invertible, which holds when $\X$ has full column rank. Then the Hessian $2\XX$ is positive definite, so this solution minimizes the SSR.}

}



\frame{
\frametitle{Numerical Example}

\textbf{Worksheet Q1}\medskip

With $n=3$, regress $y$ on an intercept and $x$ where
\[
\y=\begin{bmatrix}1\\2\\3\end{bmatrix},\quad 
\X=\begin{bmatrix}1&1\\1&3\\1&2\end{bmatrix}.
\]
\begin{enumerate}
    \item Compute $\hat{\boldsymbol{\beta}}$, fitted values $\hat\y$, residuals $\hat\u$.
    \item Verify the normal equations $\X'\hat\u=\mb{0}$.
    \item Compute $SSR=\hat\u'\hat\u$ and $\hat\sigma^2=SSR/(n-(k+1))$.
\end{enumerate}

}

\subsection{OLS Assumptions in Matrix Notation}

\frame{
\frametitle{Regression Assumptions in Matrix Notation}

\begin{itemize}
\uncover<2->{  \item \alert{I.~Linearity:} The population relationship between $\X$ and $\y$ is correctly specified and linear in parameter, $\y = \X \b +\mb{u}$\smallskip  }
\uncover<3->{  \item \alert{II.~Random Sampling}: The observed data represent a random
sample from the population described by the model.\smallskip}
\uncover<4->{  \item \alert{III.~No Perfect Collinearity}: $\X$ is a $n \times (k+1)$ matrix with rank $k+1$ (full column rank).\smallskip}
\uncover<5->{  \item \alert{IV.~Zero Conditional Mean}: $E[\mb{u}|\X]=\mb{0}$, which implies $E[\mb{u}]=\mb{0}$ and that each $u_i$ is uncorrelated with every element of $\X$.\smallskip}
\uncover<6->{  \item \alert{V.~Homoskedasticity}: $V[\mb{u|\X}]=\sig \mb{I}_{n}$ the conditional variance of the error term is constant for any value of $\X$.\smallskip}
\uncover<7->{  \item \alert{VI.~Normality}: $\mb{u}|\X \sim \mathcal{N}(\mb{0},\sig \mb{I}_{n})$ which implies that $\y|\X \sim \mathcal{N}(\X\b,\sig\mb{I}_n)$.}
\end{itemize}
}



\subsection{Properties of OLS Estimator}

\frame{
\frametitle{Expectation of a Random Vector}
The concept of expectation naturally extends to vectors of random variables:

\begin{definition}[expected value of random vector]
For a $(n \times 1)$ random vector $\mb{u}$, the expected value of $\mb{u}$, denoted $E[\mb{u}]$, is the vector of expected values:
\[
E[\mb{u}] = \left[ \begin{array}{cccc}
E[u_1] & E[u_2] & \dots & E[u_n]
 \end{array} \right]'
\]
\end{definition}
\pause
\medskip
For non-random $\mb{A}$ and $\mb{B}$, the expectation operator works as in the scalar case:
$$ E[\mb{A}\mb{u} + \mb{B}] \ = \ \mb{A}E[\mb{u}] + \mb{B} $$
If $\mb{A}$ and $\mb{B}$ are functions of $\X$, the same holds conditional on $\X$: $E[\mb{A}\mb{u} + \mb{B} \mid \X] = \mb{A}E[\mb{u} \mid \X] + \mb{B}$.


}

\frame{
\frametitle{Unbiasedness of $\hat{\b}$}

Is $E[\hat{\b}]=\b$?
\begin{align*}
\uncover<2->{\hat{\b} &= \XXinv \X'\y\, \mbox{ (linearity and no collinearity)} \\}
\uncover<3->{\hat{\b} &= \XXinv \X'(\X \b + \mb{u}) \\}
\uncover<4->{\hat{\b} &= \XXinv \XX \b + \XXinv \X' \mb{u}  \\}
\uncover<5->{\hat{\b} &= \mb{I}\b + \XXinv \X' \mb{u}\\}
\uncover<6->{\hat{\b}    &= \b + \XXinv \X'\mb{u}\,\, \\}
\uncover<7->{E[\hat{\b}|\X] &=  E[\b|\X] + E[\XXinv \X'\mb{u}|\X]\\ }
\uncover<8->{E[\hat{\b}|\X] &=  \b + \XXinv \X'E[\mb{u}|\X]\\ }
\uncover<9->{E[\hat{\b}|\X] &=  \b  \mbox{ (zero conditional mean)} }
\end{align*}
\uncover<10->{So, yes: by the law of iterated expectations, $E[\hat{\b}] = E\big[E[\hat{\b}|\X]\big] = \b$.}
}

\frame{
\frametitle{Variance-Covariance Matrix of Random Vectors}

Let's unpack the \alert{homoskedasticity} assumption $V[\mb{u|\X}]=\sig \mb{I}_{n}$.

\begin{definition}[variance-covariance matrix]
For a $(n \times 1)$ random vector $\mb{u}=\left[ \begin{array}{cccc}
u_1 & u_2 & \dots & u_n \\
 \end{array} \right]'$, its \alert{variance-covariance matrix}, denoted $V[\mb{u}]$ or also $\Sigma_{\mb{u}}$, is defined as:
\[
V[\mb{u}] = \Sigma_{\mb{u}} = \left[ \begin{array}{cccc}
\sigma^2_1 & \sigma_{12} & \dots & \sigma_{1n} \\
\sigma_{21} & \sigma^2_{2} & \dots & \sigma_{2n} \\
 \vdots & & \dots & \\
\sigma_{n1} & \sigma_{n2} & \dots & \sigma^2_{n}
 \end{array} \right]
\]
where $\sigma^2_j=V[u_j]$ and $\sigma_{ij}=Cov[u_i,u_j]$.
\end{definition}
Notice that this matrix is always symmetric.

%$$
%Var(\u) = E[\u\u'] =  \sig \mb{I}= \left[ \begin{array}{ccccc} \sigma^2 & 0 & 0 & \dots & 0 \\
%0 & \sigma^2 & 0 & \dots & 0 \\
% & & & \vdots & \\
% 0 & 0 & 0 & \dots & \sigma^2 \end{array} \right]
%$$

}

\frame{
\frametitle{Homoskedasticity in Matrix Notation}

If $V[u_i]=\sig$ for all $i=1,...,n$ and the units are independent then $V[\mb{u}] = \sig \mb{I}_n$. \\
\pause
More visually:

$$
V[\mb{u}] =  \sig \mb{I}_n=
\left[
\begin{array}{ccccc}
\sigma^2 & 0 & 0 & \dots & 0 \\
 0 & \sigma^2 & 0 & \dots & 0 \\
 & & & \vdots & \\
 0 & 0 & 0 & \dots & \sigma^2
 \end{array}
\right]
$$
\medskip
\pause
So homoskedasticity $V[\mb{u|\X}]=\sig \mb{I}_{n}$ implies that:
\begin{enumerate}
  \item $V[u_i|\X]=\sig$ for all $i$ (the variance of the errors $u_i$ does not depend on $\X$ and is constant across observations)
  \item $Cov[u_i,u_j|\X]=0$ for all $i \neq j$ (the errors are uncorrelated across observations). This holds under our random sampling assumption.
\end{enumerate}

}

%\frame{
%\frametitle{Estimation of the Error Variance}
%\small
%Given our vector of regression error terms $\mb{u}$, what is $E[\mb{u}\mb{u}']$?
%\[
%E[\mb{u}\mb{u}']=\pause\left[
%\begin{array}{ccccc}
%E[u_1^2]  & E[u_1 u_2] &  & \dots & E[u_1 u_n] \\
%E[u_2 u_1]  & E[u_2^2] &  & \dots & E[u_2 u_n] \\
%\vdots & \vdots & & \ddots & \vdots\\
%E[u_n u_1]  & E[u_n u_2] &  & \dots & E[u_n^2]
% \end{array}
%\right] = \pause \left[
%\begin{array}{ccccc}
%\sig  & 0 &  & \dots & 0 \\
%0  & \sig &  & \dots & 0 \\
%\vdots & \vdots & & \ddots & \vdots\\
%0  & 0 &  & \dots & \sig
% \end{array}
%\right]
%\]
%Recall $E[u_i]=0$ for all $i$. So $V[u_i]=E[u_i^2]-(E[u_i])^2=E[u_i^2]$ and by independence $E[u_i u_j]=E[u_i] \cdot E[u_j]=0$

%%$Cov[u_i,u_j] = E[u_i u_j] - E[u_i] \cdot E[u_j]$.% In general form, the variance-covariance matrix can be written as: $\Sigma_{u}=E[(\mb{u}-E[\mb{u}])(\mb{u}-E[\mb{u}])']$

%$$
%Var(\u) = E[\u\u'] =  \sig \mb{I}= \left[ \begin{array}{ccccc} \sigma^2 & 0 & 0 & \dots & 0 \\
%0 & \sigma^2 & 0 & \dots & 0 \\
% & & & \vdots & \\
% 0 & 0 & 0 & \dots & \sigma^2 \end{array} \right]
%$$

%}

\frame{
\frametitle{Variance of Linear Function of Random Vector}
\small
\begin{definition}[Variance of Linear Transformation of Random Vector]
Recall that for a linear transformation of a random variable $X$ we have $V[aX+b]=a^2V[X]$ with constants $a$ and $b$.\\\bigskip

There is an analogous rule for linear functions of random vectors. Let $v(\mb{u})=\mb{A}\mb{u}+\mb{B}$ be a linear transformation of a random vector $\mb{u}$ with non-random vectors or matrices $\mb{A}$ and $\mb{B}$. Then the variance of the transformation is given by:
\[
V[v(\mb{u})]=V[\mb{A}\mb{u}+\mb{B}]=\mb{A} V[\mb{u}]\mb{A}'=\mb{A} \Sigma_{\mb{u}}\mb{A}'
\]
If $\mb{A}$ and $\mb{B}$ are functions of $\X$, the same holds conditional on $\X$: $V[\mb{A}\mb{u}+\mb{B} \mid \X]=\mb{A} V[\mb{u} \mid \X]\mb{A}'$.
\end{definition}
%$$
%Var(\u) = E[\u\u'] =  \sig \mb{I}= \left[ \begin{array}{ccccc} \sigma^2 & 0 & 0 & \dots & 0 \\
%0 & \sigma^2 & 0 & \dots & 0 \\
% & & & \vdots & \\
% 0 & 0 & 0 & \dots & \sigma^2 \end{array} \right]
%$$

}

\frame{
\frametitle{Conditional Variance of $\hat{\b}$}

\textbf{Worksheet Q2}\medskip

Under homoskedasticity and independent sampling, $V[\u|\X]=\sigma^2\I_n$.  
\begin{enumerate}\itemsep1em
    \item Prove that $V[\hat{\boldsymbol{\beta}}|\X]=\sigma^2\XXinv$.  
    \item For any fixed $(k\!+\!1)\times 1$ vector $\a$, \emph{derive} $V[\a'\hat{\boldsymbol{\beta}}|\X]$ and the plug-in standard error.  
    \item Return to the data in Worksheet Q1 and report the estimated variance–covariance matrix $\widehat{V}[\hat{\boldsymbol{\beta}}]=\hat\sigma^2\XXinv$.
\end{enumerate}\bigskip

\textbf{Hint:} To estimate $V[\hat{\b}|\X]$, we replace $\sigma^2$ with its unbiased estimator
$\hat\sigma^2$, which is now written using matrix notation as:
$$ \hat\sigma^2 \ = \ \frac{SSR}{n-(k+1)} \ = \ \frac{\hat{\mb{u}}'\hat{\mb{u}}}{n-(k+1)} $$
}

% \frame{
% \frametitle{Conditional Variance of $\hat{\b}$}
% \small
% $\hat{\b}      = \b + \XXinv \X' \mb{u}$ and $E[ \hat{\b}|\X] = \b + E[ \XXinv \X' \mb{u}|\X] = \b$
% so the OLS estimator is a linear function of the errors. Thus:
% \begin{small}
% \begin{eqnarray*}
% \uncover<2->{ V[\hat{\b}|\X] &=& V[\b|\X] + V[\XXinv \X' \mb{u}|\X] \\  }
% \uncover<3->{                    &=& V[ \XXinv \X'  \mb{u}|\X]  \\  }
% \uncover<4->{                    &=&  \XXinv \X'  \alert{V[\mb{u}|\X]}  (\XXinv \X')' \,\,
% \mbox{ ($\X$ is nonrandom given $\X$})   \\ }
% \uncover<5->{                    &=&  \XXinv \X'  \alert{V[\mb{u}|\X]}  \X \XXinv   \\ }
% \uncover<6->{                    &=&  \XXinv \X'  \alert{\sig \mb{I}}  \X \XXinv \mbox{ (by homoskedasticity})\\  }
% \uncover<7->{                    &=&  \alert{\sig \mb{I}} \XXinv \X'    \X \XXinv \\  }
% \uncover<8->{                    &=& \sig  \XXinv \\  }
% \end{eqnarray*}

% \vspace{-\baselineskip}
% \uncover<9->{
% This gives the $(k+1) \times (k+1)$ \alert{variance-covariance matrix} of $\hat{\b}$.
% }
%\end{small}

% $V[\hat{\boldsymbol{\beta}}|\X]=\sigma^2\XXinv$ gives the $(k+1) \times (k+1)$ \alert{variance-covariance matrix} of $\hat{\b}$.

% To estimate $V[\hat{\b}|\X]$, we replace $\sigma^2$ with its unbiased estimator
% $\hat\sigma^2$, which is now written using matrix notation as:
% $$ \hat\sigma^2 \ = \ \frac{SSR}{n-(k+1)} \ = \ \frac{\hat{\mb{u}}'\hat{\mb{u}}}{n-(k+1)} $$
% }

%
%
%\frame{
%\frametitle{Conditional Variance of $\hat{\b}$}
%\small
%$\hat{\b}      = \b + \XXinv \X' \mb{u}$ and $E[ \hat{\b}|\X] = \b + E[ \XXinv \X' \mb{u}|\X] = \b$
%so the LS estimator is a linear function of the errors. Thus:
%\begin{small}
%\begin{align*}
%\uncover<2->{ V[\hat{\b}|\X] &= V[\b|\X] + V[\XXinv \X' \mb{u}|\X] \\  }
%\uncover<3->{                    &= V[ \XXinv \X'  \mb{u}|\X] \,\, (\mbox{ recall } V[X]=E[X^2] - (E[X])^2)   \\  }
%\uncover<4->{                    &= E[\XXinv \X'  \mb{u} \mb{u}' \X \XXinv|\X ]-\left(E[\XXinv \X'  \mb{u}|\X]E[\XXinv \X'  \mb{u}|\X]' \right) \\  }
%\uncover<5->{                    &=  \XXinv \X'  \alert{E[\mb{u} \mb{u}'|X]}  \X \XXinv-\mb{0} } \\
%\uncover<7->{                    &=  \XXinv \X'  \alert{V[\mb{u}|\X]}  \X \XXinv   \\ }
%\uncover<8->{                    &=  \XXinv \X'  \alert{\sig \mb{I}}  \X \XXinv (\mbox{by homoskedasticity})\\  }
%\uncover<9->{                    &= \sig  \XXinv \\  }
%\end{align*}
%
%\uncover<10->{The estimate of this, $\widehat{\sigma}^2  \XXinv$, is our $(k+1 \times k+1)$ \alert{variance-covariance matrix}.}
%\end{small}
%}





\frame{
\frametitle{Variance-covariance matrix of $\hat{\b}$}
The \alert{variance-covariance matrix} of the OLS estimators is given by:
$$
V[\hat{\b}|\X]= \sigma^2  \XXinv =
$$

\begin{table}
\begin{center}
\begin{tabular}{c|ccccc}
           &    $\widehat{\beta}_0$ &    $\widehat{\beta}_1$ &    $\widehat{\beta}_2$ &  $\cdots$   &    $\widehat{\beta}_k$ \\
\hline
   $\widehat{\beta}_0$ & $V[\widehat{\beta}_0]$ & $\Cov[\widehat{\beta}_0,\widehat{\beta}_1]$ &
   $\Cov[\widehat{\beta}_0,\widehat{\beta}_2]$ & $\cdots$ &  $\Cov[\widehat{\beta}_0,\widehat{\beta}_k]$   \\

   $\widehat{\beta}_1$ & $\Cov[\widehat{\beta}_0,\widehat{\beta}_1]$ & $V[\widehat{\beta}_1]$ &
    $\Cov[\widehat{\beta}_1,\widehat{\beta}_2]$ &  $\cdots$    &  $\Cov[\widehat{\beta}_1,\widehat{\beta}_k]$ \\

   $\widehat{\beta}_2$ & $\Cov[\widehat{\beta}_0,\widehat{\beta}_2]$ & $\Cov[\widehat{\beta}_1,\widehat{\beta}_2]$ & $V[\widehat{\beta}_2]$ &   $\cdots$  &   $\Cov[\widehat{\beta}_2,\widehat{\beta}_k]$  \\
    $\vdots$    &  $\vdots$  & $\vdots$    &  $\vdots$    &    $\ddots$        &    $\vdots$        \\
   $\widehat{\beta}_k$ & $\Cov[\widehat{\beta}_0,\widehat{\beta}_k]$ & $\Cov[\widehat{\beta}_k,\widehat{\beta}_1]$ & $\Cov[\widehat{\beta}_k,\widehat{\beta}_2]$ &     $\cdots$       & $V[\widehat{\beta}_k]$ \\
\end{tabular}
\end{center}
\end{table}
}

\begin{frame}[fragile]
\frametitle{Variance-covariance matrix of $\hat{\b}$}

\spacingset{0.8}
% Output of code/fig_chile_matrices.R (saved in code/fig_chile_matrices.txt).
\begin{Verbatim}[fontsize=\footnotesize, frame=single, label=R Code, commandchars=\\\{\}]
> fit <- lm(vote1 ~ fem + educ + age, data = d)
> summary(fit)
Coefficients:
              Estimate Std. Error t value Pr(>|t|)
(Intercept)  0.4917751  \alert{0.0525717}   9.354  < 2e-16 ***
fem          0.1284292  \alert{0.0236221}   5.437 6.21e-08 ***
educ        -0.1051514  \alert{0.0164543}  -6.391 2.13e-10 ***
age          0.0033573  \alert{0.0008235}   4.077 4.77e-05 ***
---
\end{Verbatim}

\footnotesize
We can pull out the variance-covariance matrix $\hat\sigma^2 (\X'\X)^{-1}$ in \texttt{R} from the \texttt{lm()} object:
\begin{Verbatim}[fontsize=\footnotesize, frame=single, label=R Code, commandchars=\\\{\}]
> V <- vcov(fit)
> V
              (Intercept)           fem          educ           age
(Intercept)  2.763788e-03 -3.717198e-04 -6.543766e-04 -3.324904e-05
fem         -3.717198e-04  5.580053e-04  3.781282e-05  9.631845e-07
educ        -6.543766e-04  3.781282e-05  2.707446e-04  3.860967e-06
age         -3.324904e-05  9.631845e-07  3.860967e-06  6.780985e-07

> sqrt(diag(V))
 (Intercept)          fem         educ          age
\alert{0.0525717389 0.0236221360 0.0164543200 0.0008234674}
\end{Verbatim}
These are the conventional SEs. They assume homoskedasticity (Assumption V), which a 0/1 outcome makes doubtful. Here robust SEs (Lecture 6.3) are almost the same: 0.0237 instead of 0.0236 for \texttt{fem}.

\end{frame}

\begin{frame}
\frametitle{Consistency and Asymptotic Normality}
\textbf{Worksheet Q3}\medskip

Suppose 
\begin{itemize}
\item the true model is $y_i = \mb{x}_i\boldsymbol{\beta} + u_i$, where $\mb{x}_i$ is row $i$ of $\X$ (a $1 \times (k+1)$ row vector);
\item with $\E[u_i|\mb{x}_i]=0$ and $\E[u_i^2]<\infty$;
\item the sample $\{(y_i,\mb{x}_i)\}_{i=1}^n$ is i.i.d.; 
\item $\E[\mb{x}_i'\mb{x}_i]$ is finite and nonsingular.  
\end{itemize}\medskip
 
\begin{enumerate}\itemsep1em
    \item Show that $\hat{\boldsymbol{\beta}}$ is a consistent estimator of $\boldsymbol{\beta}$.
    \item Assume further that $E[u_i^2|\mb{x}_i]=\sigma^2<\infty$, derive the asymptotic distribution of $\hat{\boldsymbol{\beta}}$.
\end{enumerate}


\end{frame}


% \begin{frame}
% \frametitle{Consistency of $\hat{\b}$}
% \small
% To show consistency, we rewrite the OLS estimator in terms of sample means so that we can
% apply LLN.\\
% \medskip
% \pause
% First, note that a matrix cross product can be written as a sum of vector products:
% $$ \XX \ = \ \sum_{i=1}^n \mb{x}_i'\mb{x}_i \quad \text{and} \quad
%    \X'\y \ = \ \sum_{i=1}^n \mb{x}_i'y_i $$
% where $\mb{x}_i$ is the $1 \times (k+1)$ \alert{row} vector of predictor values for unit $i$.\\
% \pause
% \medskip
% Now we can rewrite the OLS estimator as,
% \begin{align*}
% \hat{\b} &= \left( \sum\nolimits_{i=1}^n \mb{x}_i'\mb{x}_i   \right)^{-1} \left( \sum\nolimits_{i=1}^n \mb{x}_i' y_i \right) \\
%  &= \left( \sum\nolimits_{i=1}^n \mb{x}_i'\mb{x}_i   \right)^{-1} \left( \sum\nolimits_{i=1}^n \mb{x}_i'(\mb{x}_i\b + u_i) \right) \\
%  &= \b +  \left( \sum\nolimits_{i=1}^n \mb{x}_i'\mb{x}_i   \right)^{-1} \left( \sum\nolimits_{i=1}^n \mb{x}_i' u_i \right) \\
%  &= \b +  \left(\frac{1}{n} \sum\nolimits_{i=1}^n \mb{x}_i'\mb{x}_i   \right)^{-1} \left(\frac{1}{n} \sum\nolimits_{i=1}^n \mb{x}_i' u_i \right)
% \end{align*}
% \end{frame}

% \begin{frame}
% \frametitle{Consistency of $\hat{\b}$}
% \small
% Now let's apply the LLN to the sample means:
% \begin{eqnarray*}
% & & \left(\frac{1}{n} \sum_{i=1}^n \mb{x}_i'\mb{x}_i   \right) \ \stackrel{p}{\longrightarrow} \ E[\mb{x}_i'\mb{x}_i], \ \text{a $(k+1)\times(k+1)$ nonsingular matrix.} \\
% & & \left(\frac{1}{n} \sum_{i=1}^n \mb{x}_i' u_i \right) \ \stackrel{p}{\longrightarrow} \ E[\mb{x}_i' u_i] \ = \ 0, \ \text{by the zero cond. mean assumption.}
% \end{eqnarray*}
% \pause

% Therefore, we have
% \begin{eqnarray*}
% \plim ( \hat{\b})  & = & \b + \left(E[\mb{x}_i'\mb{x}_i]\right)^{-1} \cdot 0 \\
%   & = & \b.
% \end{eqnarray*}

% \pause
% \medskip
% We can also show the asymptotic normality of $\hat{\b}$ using a similar argument but
% with the CLT (see Wooldrigde Theorem 11.2 for details).

% \end{frame}



\frame{
\frametitle{Properties of the OLS Estimator: Summary}
\small
\begin{Theorem}
 Under Assumptions I--VI, the $(k+1) \times 1$ vector of OLS estimators $\hat{\b}$, conditional on $\X$,  follows a \alert{multivariate normal distribution} with mean $\b$ and variance-covariance matrix $\sig  \XXinv$:
$$
\hat{\b}|\X \sim \mathcal{N}\left(\b,\sig \XXinv \right)
$$
\vspace{-\baselineskip}
\begin{itemize}
\addtolength{\itemsep}{-.25\baselineskip}
  \item Each element of $\hat{\b}$ (i.e.\ $\hat{\beta}_0,\dots,\hat{\beta}_k$) is normally distributed, and $\hat{\b}$ is an unbiased estimator of $\b$ as $E[\hat{\b}]=\b$
 \item Variances and covariances are given by $V[\hat{\b}|\X]=\sig  \XXinv$
  \item An unbiased estimator for the error variance $\sig$ is given by $$\widehat{\sigma}^2=\frac{\hat{\mb{u}}'\hat{\mb{u}}}{n-(k+1)}$$
\vspace{-\baselineskip}
  \item With a large sample, $\hat{\b}$ approximately follows the same distribution
    under Assumptions I--V only, i.e., without assuming the normality of $\mb{u}$.
 % \item $\hat{\b}$ is also the best linear unbiased estimator
\end{itemize}
\end{Theorem}
}



%\frame{
%\frametitle{Unconditional Variance of $\hat{\b}$ }
%
%The preceding variance analysis was done conditional on the observed data. We can obtain the unconditional variance by averaging over $\X$ (the conditional distribution).\\\bigskip
%
%Recall the variance decomposition. For any RVs $Z$ and $X$ we have:$$
%V[Z]=V[E[Z|W]]+E[V[Z|W]]
%$$
%So the unconditional variance is obtained by:
%\begin{small}
%\begin{align*}
%\uncover<2->{ V[\hat{\b}] &= E_{X}[V[\hat{\b}|\X]]+V_{X}[E[\hat{\b}|X]]\\  }
%\uncover<3->{  &= E_{X}[V[\hat{\b}|\X]]+V_{X}[\b]\\  }
%\uncover<4->{                    &= E_{X}[V[\hat{\b}|\X]]   \\  }
%\uncover<5->{                    &= E_{X}[ \sig  \XXinv  ] \\  }
%\uncover<6->{                    &= \sig E_{X}[   \XXinv  ]}
%\end{align*}
%\uncover<6->{For fixed $\X$ or asymptotically this is the same as the conditional variance.}
%\end{small}
%}

\section{Frisch–Waugh–Lovell Theorem}

\subsection{Projection and Annihilator Matrices}


\frame{
\frametitle{Review: Regression in Matrix Algebra}
\textbf{Worksheet Q4}\medskip\small

Assume the classical linear regression model satisfying Assumptions I–V:
\[
\mathbf{y} = \mathbf{X}\boldsymbol{\beta} + \mathbf{u},
\]
where $\mathbf{X}$ is an $n \times (k+1)$ matrix of full column rank, and $\mathbf{y}$ is an $n \times 1$ vector of responses.  

\begin{enumerate}
\item What does ${\bf (X'X)^{-1}X'y}$ represent in the context of OLS estimation?  

\item What is the interpretation of ${\bf X(X'X)^{-1}X'y}$?  

\item What does ${\bf y} - {\bf X(X'X)^{-1}X'y}$ represent?  

\item What does $[{\bf y} - {\bf X(X'X)^{-1}X'y}]'[{\bf y} - {\bf X(X'X)^{-1}X'y}]$ compute?  

\item What does $[{\bf y} - {\bf X(X'X)^{-1}X'y}]'[{\bf y} - {\bf X(X'X)^{-1}X'y}]/(n-(k+1))$ estimate?  

\item What is the interpretation of 
\[
{\bf (X'X)^{-1}}\,\frac{[{\bf y} - {\bf X(X'X)^{-1}X'y}]'[{\bf y} - {\bf X(X'X)^{-1}X'y}]}{(n-(k+1))}\,?
\]  
\end{enumerate}
}



\frame{
\frametitle{Projection Matrix and Annihilator Matrix}
\small
\vspace{-\baselineskip}
\begin{eqnarray*}
\hat{\u} & = & \y - \hat{\y} =  \y - \X \hat{\b} \\
 & = & \y - \alert{\X \XXinv \X'} \y \equiv \y - \alert{\mb{P}}\y \\ \pause
 & = & \blue{(\mb{I} - \mb{P})}\y  \equiv  \blue{\mb{M}}\y
\end{eqnarray*}
\pause
$\mb{P}$ is called the \alert{projection matrix}:
\begin{itemize}
\item The above shows $\hat{\y} = \mb{P}\y$, i.e, premultiplying a vector
by $\mb{P}$ turns the vector into its linear fit with respect to $\X$.
\item $\mb{P}$ is therefore also known as the \alert{hat matrix} 
\end{itemize}
\pause
$\mb{M}$ is called the \alert{annihilator matrix} (because $\mb{MX=0}$) and more intuitively the ``residual maker'' matrix:
\begin{itemize}
\item The above shows $\hat{\u} = \mb{M}\y$, i.e., pre-multiplying a vector
by $\mb{M}$ turns the vector into the residual from its regression on $\X$.
\item $\mb{M}$ is therefore a ``residual creator'' matrix
\end{itemize}
\pause
$\mb{P}$ and $\mb{M}$ are both symmetric and \alert{idempotent} (i.e.
$\mb{P}\mb{P}=\mb{P}$ and $\mb{M}\mb{M}=\mb{M}$).
}

\begin{frame}
\frametitle{The Geometry of Least Squares}
% New slide (2026), drawn in TikZ. Paid for by merging the two FWL slides below.
\begin{columns}[c]
\begin{column}{0.55\textwidth}
\centering
\begin{tikzpicture}[x=0.8cm, y=0.8cm, >=Stealth, font=\footnotesize]
  % the column space of X: a plane, drawn in perspective
  \fill[CourseMist] (0,0) -- (6.6,0) -- (8.2,2.6) -- (1.6,2.6) -- cycle;
  \draw[gray!60] (0,0) -- (6.6,0) -- (8.2,2.6) -- (1.6,2.6) -- cycle;
  \node[gray, anchor=south east, font=\scriptsize] at (6.1,0.05) {column space of $\X$};
  % the two columns of X (n = 3: an intercept and one regressor)
  \draw[->, gray, thick] (1.1,0.55) -- (2.6,0.3) node[below, text=CourseInk] {$\mb{1}$};
  \draw[->, gray, thick] (1.1,0.55) -- (1.55,2.0) node[left, text=CourseInk] {$\x$};
  % y, and a fit X beta-tilde that is not the closest point
  \draw[->, CourseInk, very thick] (1.1,0.55) -- (4.2,5.0) node[pos=0.5, above left] {$\y$};
  \draw[gray, dashed] (6.9,1.0) -- (4.2,5.0);
  \fill[gray] (6.9,1.0) circle (1.8pt) node[below right, text=CourseInk] {$\X\tilde{\b}$};
  % the OLS fit and the residual
  \uncover<2->{
    \draw[->, Cardinal, very thick] (1.1,0.55) -- (4.2,1.35)
      node[pos=0.62, below right, text=Cardinal] {$\hat{\y} = \P\y$};
    \draw[->, CourseBlue, very thick] (4.2,1.35) -- (4.2,5.0)
      node[pos=0.3, left, text=CourseBlue] {$\hat{\u} = \M\y$};
    % right-angle mark at y-hat: 0.5 units toward the origin, then 0.5 up
    % (large enough that the arrowhead of y-hat does not cover its corner)
    \draw[CourseInk, thin] (3.716,1.225) -- (3.716,1.725) -- (4.2,1.85);
    \fill[Cardinal] (4.2,1.35) circle (1.8pt);
  }
\end{tikzpicture}
\end{column}
\begin{column}{0.43\textwidth}
\small
With $n = 3$ observations, $\y$, $\mb{1}$ and $\x$ are vectors in 3-D space.
\begin{itemize}
\item Every fit $\X\tilde{\b}$ lies in the plane spanned by the columns of $\X$: the \alert{column space} of $\X$.
\item<2-> OLS picks the point of the plane closest to $\y$: $\hat{\y} = \X\hat{\b} = \P\y$.
\item<2-> The shortest path from $\y$ to the plane meets it at a right angle. So $\hat{\u} = \M\y$ is orthogonal to every column of $\X$: $\X'\hat{\u} = \mb{0}$.
\item<2-> $\y = \P\y + \M\y$: $\P$ keeps the part of $\y$ in the plane; $\M$ keeps the part perpendicular to it.
\end{itemize}
\end{column}
\end{columns}
\slidesource{Adapted from Blackwell, \emph{A User's Guide to Statistical Inference and Regression}, ch.~7}
\end{frame}




\frame{
\frametitle{Projection Matrix and Annihilator Matrix}
\textbf{Worksheet Q5}\medskip\small

Assume the classical linear regression model satisfying Assumptions I–V: $\mathbf{y} = \mathbf{X}\boldsymbol{\beta} + \mathbf{u}$, where $\mathbf{X}$ is an $n \times (k+1)$ matrix of full column rank, and $\mathbf{y}$ is an $n \times 1$ vector of responses.  Define: 
\[
\mathbf{Q}=\X'\X,\qquad \mathbf{A}=\mathbf{Q}^{-1}\X',\qquad \P=\X\mathbf{A},\qquad \M=\I-\P.
\]
\vspace{-2em}\begin{enumerate}
\item Express the residuals from regressing $\y$ on $\X$ as a function of $\M$.  

\item Show that $\P$ is symmetric and idempotent. 

\item Show that $\M$ is idempotent (verify $\M'\M=\M$ and $\M^2=\M$).  

\item Show that $\mathbf{A}\P=\mathbf{A}$, $\mathbf{A}\M=\mathbf{0}$, $\M\P=\mathbf{0}$, and $\P\M=\mathbf{0}.$

\item Show that $\P\X=\X$, $\M\X=\mathbf{0}$.

\item Show that $\P\,\widehat{\y}=\widehat{\y}$, $\P\,\widehat{\mathbf{u}}=\mathbf{0}.$  

\item Show that $\M\,\widehat{\y}=\mathbf{0}$, $\M\,\widehat{\mathbf{u}}=\widehat{\mathbf{u}}$.  

\item Show that $\X'\widehat{\y}=\X'\y$.  

\item Show that $\widehat{\mathbf{u}}'\widehat{\mathbf{u}}=\y'\M\y=\y'\y-\widehat{\y}'\widehat{\y}$.  
\end{enumerate}
}



\subsection{Frisch–Waugh–Lovell Theorem}

\frame{
\frametitle{Frisch–Waugh–Lovell Theorem}
% 2026: old slides 41 and 42 merged into this one (both originals are in
% ../_backup/06-1MultipleRegressionMechanics_removed-2026-09-17.tex).
\vspace{-\baselineskip}
\begin{itemize}
\item In multiple regression
\[
\y = \X_1\b_1 + \X_2\b_2 + \u,
\]
we often want to interpret the coefficient $\b_2$ as the effect of $\X_2$ \emph{controlling for} $\X_1$.
\item The FWL theorem shows that we can obtain $\hat{\b}_2$ in two equivalent ways:
\begin{enumerate}
\item Regress $\y$ directly on both $\X_1$ and $\X_2$ (full regression), or
\item First remove from $\y$ and $\X_2$ the parts explained by $\X_1$ (``partial out'' $\X_1$), then regress the residuals of $\y$ on those of $\X_2$.
\end{enumerate}
\item The theorem provides a clean geometric interpretation using projection and residual-maker matrices.
\end{itemize}
\pause
\begin{block}{Frisch–Waugh–Lovell Theorem}
The coefficient $\hat{\b}_2$ from regressing $\y$ on $(\X_1,\X_2)$ equals that from regressing the residualized $\y$ on residualized $\X_2$:
\[
\hat{\b}_2 = (\X_2'\M_1\X_2)^{-1}\X_2'\M_1\y, \qquad \M_1 = \I - \X_1(\X_1'\X_1)^{-1}\X_1'
\]
The two regressions also leave the same residuals.
\end{block}
}


% \frame{
% \frametitle{Setup}
% Consider the linear model:
% \[
% \y = \X_1 \a + \X_2 \b + \u
% \]
% \vspace{-\baselineskip}
% \begin{itemize}
% \item $\X_1$ and $\X_2$ are full-column-rank matrices.
% \item Goal: estimate $\b$ net of $\X_1$ (partialling out $\X_1$).
% \item OLS normal equations (FOCs):
% \[
% \X_1'(\y-\X_1\hat{\a}-\X_2\hat{\b})=0, \quad
% \X_2'(\y-\X_1\hat{\a}-\X_2\hat{\b})=0.
% \]
% \end{itemize}
% }

% \frame{
% \frametitle{Step 1: Eliminate $\hat{\a}$ Using Normal Equations}
% \small
% From the first equation:
% \[
% \hat{\a} = (\X_1'\X_1)^{-1}\X_1'(\y-\X_2\hat{\b})
% \]
% Substitute into the second:
% \[
% \X_2'[\I - \X_1(\X_1'\X_1)^{-1}\X_1'](\y-\X_2\hat{\b})=0
% \]
% \pause
% Define the residual-maker matrix for $\X_1$:
% \[
% \M_1 = \I - \X_1(\X_1'\X_1)^{-1}\X_1'
% \]
% Then
% \[
% \X_2'\M_1\y = \X_2'\M_1\X_2\hat{\b}.
% \]
% }

% \frame{
% \frametitle{Step 2: The Partial Regression Formula}
% Solve for $\hat{\b}$:
% \[
% \hat{\b} = (\X_2'\M_1\X_2)^{-1}\X_2'\M_1\y
% \]
% \pause
% Interpretation:
% \begin{itemize}
% \item $\M_1\y$ are residuals from regressing $\y$ on $\X_1$.
% \item $\M_1\X_2$ are residuals from regressing each column of $\X_2$ on $\X_1$.
% \item Hence, $\hat{\b}$ equals the coefficient from regressing $\M_1\y$ on $\M_1\X_2$.
% \end{itemize}
% \[
% \hat{\b} = \big((\M_1\X_2)'(\M_1\X_2)\big)^{-1}(\M_1\X_2)'(\M_1\y)
% \]
% }

% \frame{
% \frametitle{Step 3a: Residuals and the Residual-Maker Matrix}
% \small

% We begin with the full regression model:
% \[
% \y = \X_1\boldsymbol{\alpha} + \X_2\boldsymbol{\beta} + \u.
% \]

% The OLS residuals are:
% \[
% \hat{\u} = \y - \X_1\hat{\boldsymbol{\alpha}} - \X_2\hat{\boldsymbol{\beta}}.
% \]

% \pause
% Let $\M_1 = \I - \P_{X_1}$ be the \textit{residual-maker matrix} for $\X_1$.

% By construction, $\M_1$ removes any linear dependence on $\X_1$, it is also symmetric and idempotent:
% \[
% \M_1\X_1 = \mathbf{0},\qquad  \M_1' = \M_1, \qquad \M_1^2 = \M_1.
% \]
% \pause
% Now, pre-multiply the full regression by $\M_1$:
% \[
% \M_1\y = \M_1\X_1\boldsymbol{\alpha} + \M_1\X_2\boldsymbol{\beta} + \M_1\u.
% \]

% Because $\M_1\X_1 = 0$, this simplifies to:
% \[
% \M_1\y = \M_1\X_2\boldsymbol{\beta} + \M_1\u.
% \]

% This suggests that once the variation in $\y$ and $\X_2$ explained by $\X_1$ is removed, 
% the remaining regression involves only $\M_1\y$ and $\M_1\X_2$.
% }

% \frame{
% \frametitle{Step 3b: Projection Equivalence and Decomposition}
% \small

% Using the above relationship, we can re-express the OLS residuals step by step:
% \[
% \begin{aligned}
% \hat{\u}
% &= \y - \X_1\hat{\boldsymbol{\alpha}} - \X_2\hat{\boldsymbol{\beta}} \\[3pt]
% &= (\I - \P_{X_1})\y - (\I - \P_{X_1})\X_2\hat{\boldsymbol{\beta}} \\[3pt]
% &= \M_1(\y - \X_2\hat{\boldsymbol{\beta}}).
% \end{aligned}
% \]

% \pause
% Therefore, $\hat{\u}$ equals the residuals from regressing the \textit{residualized outcome} $\M_1\y$ 
% on the \textit{residualized regressor} $\M_1\X_2$.

% \pause
% This implies the equivalence:
% \[
% \text{Residuals from } (\y \sim \X_1, \X_2) 
% \;=\;
% \text{Residuals from } (\M_1\y \sim \M_1\X_2).
% \]

% \pause
% Moreover, the fitted values from the full regression decompose orthogonally:
% \[
% \P_X\y = \P_{X_1}\y + \P_{\M_1 X_2}\y,
% \]
% where:
% \[
% \P_{X_1}\y \text{ captures the variation in } \y \text{ explained by } \X_1, 
% \]
% \[
% \P_{\M_1 X_2}\y \text{ captures the additional variation explained by } \X_2 \text{ after removing } \X_1.
% \]

% \medskip
% This geometric decomposition is the essence of the Frisch–Waugh–Lovell theorem.
% }




\end{document}
